% Generated mechanically from the frozen spaces-duality.tex target.
% Do not edit this generated file; edit the bound locale source instead.

% OK, start here.
%

\setcounter{chapter}{85}
\chapter{空间的对偶性}



\phantomsection
\label{spaces-duality-section-phantom}


\section{引言}
\label{spaces-duality-section-introduction}

\noindent
本章是概形情形相应章节的类似版本，见《概形的对偶性》第
\ref{duality-section-introduction} 节。这里的论述与下列论文中的论述相近：
\cite{Neeman-Grothendieck}, \cite{LN},
\cite{Lipman-notes}、\cite{Neeman-improvement}。




\section{代数空间上的对偶复形}
\label{spaces-duality-section-dualizing-spaces}

\noindent
设 $U$ 是局部 Noether 概形。以 $\mathcal{O}_\etale$ 表示 $U$ 上的
结构层（定义在 $U$ 的小平展站点上）。若存在《概形的对偶性》第
\ref{duality-section-dualizing-schemes} 节意义下的对偶复形
$\omega_U^\bullet$，使得
$K = \epsilon^*(\omega_U^\bullet)$，则称对象
$K \in D_\QCoh(\mathcal{O}_\etale)$ 是 $U$ 上的对偶复形。
这里，$\epsilon^* : D_\QCoh(\mathcal{O}_U) \to
D_\QCoh(\mathcal{O}_\etale)$ 是《空间的导出范畴》引理
\ref{spaces-perfect-lemma-derived-quasi-coherent-small-etale-site}
中的等价。《概形的对偶性》第
\ref{duality-section-dualizing-schemes} 节中研究的
$\omega_U^\bullet$ 的大部分性质，经《空间的导出范畴》第
\ref{spaces-perfect-section-derived-quasi-coherent-etale} 与
\ref{spaces-perfect-section-spell-out} 节的讨论由 $K$ 继承。

\medskip\noindent
我们把局部 Noether 代数空间上的对偶复形定义为这样一种复形：
它在平展局部来自相应概形上的对偶复形。

\begin{lemma}
\label{spaces-duality-lemma-equivalent-definitions}
设 $S$ 是概形，$X$ 是 $S$ 上的局部 Noether 代数空间，
$K$ 是 $D_\QCoh(\mathcal{O}_X)$ 的对象。下列条件等价：
\begin{enumerate}
\item 对每个以概形 $U$ 为源的平展态射 $U \to X$，限制
$K|_U$ 都是 $U$ 上的对偶复形（含义如上）。
\item 存在以概形 $U$ 为源的满平展态射 $U \to X$，使得
$K|_U$ 是 $U$ 上的对偶复形。
\end{enumerate}
\end{lemma}

\begin{proof}
设 $U \to X$ 是以概形 $U$ 为源的满平展态射，并设
$V \to X$ 是以概形 $V$ 为源的平展态射。则
$$
U \leftarrow U \times_X V \rightarrow V
$$
是概形间的平展态射，其中指向 $V$ 的箭头是满的。因此，由
《概形的对偶性》引理 \ref{duality-lemma-descent-ascent} 可知：
若 $K|_U$ 是 $U$ 上的对偶复形，则 $K|_V$ 是 $V$ 上的对偶复形。
\end{proof}

\begin{definition}
\label{spaces-duality-definition-dualizing-scheme}
设 $S$ 是概形，$X$ 是 $S$ 上的局部 Noether 代数空间。
若 $D_\QCoh(\mathcal{O}_X)$ 的对象 $K$ 满足引理
\ref{spaces-duality-lemma-equivalent-definitions} 中的等价条件，则称 $K$ 为
{\it 对偶复形}。
\end{definition}

\begin{lemma}
\label{spaces-duality-lemma-affine-duality}
设 $A$ 是 Noether 环，$X = \Spec(A)$。以 $\mathcal{O}_\etale$
表示 $X$ 上的结构层（定义在 $X$ 的小平展站点上）。设 $K,L$ 是
$D(A)$ 的对象。
若 $K \in D_{\textit{Coh}}(A)$ 且 $L$ 具有有限内射维数，则
$$
\epsilon^*\widetilde{R\Hom_A(K, L)} =
R\SheafHom_{\mathcal{O}_\etale}(\epsilon^*\widetilde{K},
\epsilon^*\widetilde{L})
$$
在 $D(\mathcal{O}_\etale)$ 中成立，其中
$\epsilon : (X_\etale, \mathcal{O}_\etale) \to (X, \mathcal{O}_X)$
如《空间的导出范畴》第
\ref{spaces-perfect-section-derived-quasi-coherent-etale} 节所述。
\end{lemma}

\begin{proof}
由《概形的对偶性》引理 \ref{duality-lemma-affine-duality}，有典范同构
$$
\widetilde{R\Hom_A(K, L)} =
R\SheafHom_{\mathcal{O}_X}(\widetilde{K}, \widetilde{L})
$$
于 $D(\mathcal{O}_X)$ 中。又有典范映射
$$
\epsilon^*R\Hom_{\mathcal{O}_X}(\widetilde{K}, \widetilde{L})
\longrightarrow
R\SheafHom_{\mathcal{O}_\etale}(\epsilon^*\widetilde{K},
\epsilon^*\widetilde{L})
$$
于 $D(\mathcal{O}_\etale)$ 中，见《站点上的上同调》评注
\ref{sites-cohomology-remark-prepare-fancy-base-change}。
我们将证明该箭头两端具有同构的上同调层，但略去验证这一同构
确由该箭头给出。

\medskip\noindent
可以假设 $L$ 由内射 $A$-模组成的有限复形 $I^\bullet$ 给出。
对 $I^\bullet$ 的长度作归纳，并利用这些构造与特异三角的相容性，
可归约到 $L=I[0]$ 的情形，其中 $I$ 是内射 $A$-模。回忆
% P10-E0246
$R\SheafHom_{\mathcal{O}_\etale}(\epsilon^*\widetilde{K},
\epsilon^*\widetilde{L}))$
的上同调层是下述预层的层化：它把 $X$ 上平展的 $U$ 映到
$\epsilon^*\widetilde{K}$ 与 $\epsilon^*\widetilde{L}$ 在
$U_\etale$ 上的限制之间的第 $i$ 个 Ext 群。见《站点上的上同调》
引理 \ref{sites-cohomology-lemma-section-RHom-over-U}。
若 $U=\Spec(B)$ 是仿射的，则由《空间的导出范畴》引理
\ref{spaces-perfect-lemma-derived-quasi-coherent-small-etale-site} 及
《概形的导出范畴》引理
\ref{perfect-lemma-affine-compare-bounded}
给出的等价，该 Ext 群等于
$\text{Ext}^i_B(K \otimes_A B, L \otimes_A B)$（这里还使用了
《空间的导出范畴》评注
\ref{spaces-perfect-remark-match-total-direct-images} 中详述的相容性）。
由于 $A\to B$ 是平展的，由《对偶复形》引理
\ref{dualizing-lemma-injective-goes-up} 可知
$I\otimes_A B$ 是内射 $B$-模。因此
\begin{align*}
\Ext^n_B(K \otimes_A B, I \otimes_A B)
& =
\Hom_B(H^{-n}(K \otimes_A B), I \otimes_A B) \\
& =
% P10-E0247
\Hom_{A_f}(H^{-n}(K) \otimes_A B, I \otimes_A B) \\
& =
\Hom_A(H^{-n}(K), I) \otimes_A B \\
& =
\text{Ext}^n_A(K, I) \otimes_A B
\end{align*}
% P10-E0248
倒数第二个等号成立，是因为 $H^{-n}(K)$ 是有限 $A$-模；见
《代数详论》引理
\ref{more-algebra-lemma-pseudo-coherence-and-base-change-ext}。
因此，引理中等式两端的上同调层相同。
\end{proof}

\begin{lemma}
\label{spaces-duality-lemma-dualizing-spaces}
设 $S$ 是概形，$X$ 是 $S$ 上的局部 Noether 代数空间，
$K$ 是 $X$ 上的对偶复形。则 $K$ 是
$D_{\textit{Coh}}(\mathcal{O}_X)$ 的对象，并且
$D = R\SheafHom_{\mathcal{O}_X}(-, K)$ 诱导反等价
$$
D :
D_{\textit{Coh}}(\mathcal{O}_X)
\longrightarrow
D_{\textit{Coh}}(\mathcal{O}_X)
$$
，并配有典范同构 $\text{id}\to D\circ D$。若 $X$ 拟紧，则
$D$ 交换 $D^+_{\textit{Coh}}(\mathcal{O}_X)$ 与
$D^-_{\textit{Coh}}(\mathcal{O}_X)$，并诱导等价
$D^b_{\textit{Coh}}(\mathcal{O}_X) \to
D^b_{\textit{Coh}}(\mathcal{O}_X)$。
\end{lemma}

\begin{proof}
取以仿射概形 $U$ 为源的平展态射 $U\to X$。设 $U=\Spec(A)$，
并令 $\omega_A^\bullet$ 为与 $K|_U$ 对应的 $A$ 的对偶复形，
其含义如引理 \ref{spaces-duality-lemma-equivalent-definitions} 及《概形的对偶性》
引理 \ref{duality-lemma-equivalent-definitions} 所述。由引理
\ref{spaces-duality-lemma-affine-duality}，下图交换：
$$
\xymatrix{
D_{\textit{Coh}}(A) \ar[r] \ar[d]_{R\Hom_A(-, \omega_A^\bullet)} &
D_{\textit{Coh}}(\mathcal{O}_\etale)
\ar[d]^{R\SheafHom_{\mathcal{O}_\etale}(-, K|_U)} \\
D_{\textit{Coh}}(A) \ar[r] &
D(\mathcal{O}_\etale)
}
$$
其中 $\mathcal{O}_\etale$ 是 $U$ 的小平展站点上的结构层。
由于 $R\SheafHom$ 的形成与限制相交换，可知 $D$ 把
$D_{\textit{Coh}}(\mathcal{O}_X)$ 映入
$D_{\textit{Coh}}(\mathcal{O}_X)$。此外，典范映射
$$
L \longrightarrow
R\SheafHom_{\mathcal{O}_X}(R\SheafHom_{\mathcal{O}_X}(L, K), K)
$$
（《站点上的上同调》引理
\ref{sites-cohomology-lemma-internal-hom-evaluate}）
对 $D_{\textit{Coh}}(\mathcal{O}_X)$ 中每个 $L$ 都是同构；这是因为，
由《对偶复形》引理 \ref{dualizing-lemma-dualizing}，在上述每个 $U$
上这一结论都成立。拟紧情形下关于函子 $D$ 的有界性性质，也由
《对偶复形》引理 \ref{dualizing-lemma-dualizing} 中的相应结论推出。
\end{proof}

\noindent
设 $(\mathcal{C},\mathcal{O})$ 是赋环站点。回忆：若
$D(\mathcal{O})$ 的对象 $L$ 对于由导出张量积给出的
% P10-E0249
$D(\mathcal{O}_X)$ 上的对称幺半结构是可逆对象，则称它
{\it 可逆}。由《站点上的上同调》引理
\ref{sites-cohomology-lemma-invertible-derived} 可知，
% P10-E0250
这等价于 $L$ 是完美的；并且若 $(\mathcal{C},\mathcal{O})$
是局部赋环站点，则对 $\mathcal{C}$ 的每个对象 $U$，存在
$\mathcal{C}$ 中 $U$ 的覆盖 $\{U_i\to U\}$，使得对某些整数 $n_i$ 有
$L|_{U_i}\cong\mathcal{O}_{U_i}[-n_i]$。

\medskip\noindent
设 $S$ 是概形，$X$ 是 $S$ 上的代数空间。若
$D(\mathcal{O}_X)$ 中的 $L$ 可逆，则存在不交并分解
$X=\coprod_{n\in\mathbf{Z}}X_n$，使得 $L|_{X_n}$ 是位于次数 $n$
的可逆模。特别地，$L=\bigoplus H^n(L)[-n]$；这给出了一个表示 $L$
的良定义 $\mathcal{O}_X$-模复形（其微分均为零）。

\begin{lemma}
\label{spaces-duality-lemma-dualizing-unique-spaces}
设 $S$ 是概形，$X$ 是 $S$ 上的局部 Noether 代数空间。
若 $K$ 与 $K'$ 是 $X$ 上的对偶复形，则 $K'$ 同构于
$K\otimes_{\mathcal{O}_X}^\mathbf{L}L$，其中 $L$ 是
$D(\mathcal{O}_X)$ 中的某个可逆对象。
\end{lemma}

\begin{proof}
置
$$
L = R\SheafHom_{\mathcal{O}_X}(K, K')
$$
它是 $D(\mathcal{O}_X)$ 中的可逆对象，因为这一点在仿射局部成立；
可使用引理 \ref{spaces-duality-lemma-affine-duality} 以及《对偶复形》引理
\ref{dualizing-lemma-dualizing-unique} 及其证明。由同样的理由，
求值映射 $L\otimes_{\mathcal{O}_X}^\mathbf{L}K\to K'$ 是同构。
\end{proof}

\begin{lemma}
\label{spaces-duality-lemma-dimension-function-scheme}
设 $S$ 是概形，$X$ 是 $S$ 上局部 Noether 且拟分离的代数空间，
$\omega_X^\bullet$ 是 $X$ 上的对偶复形。则对 $X$，函数
% P10-E0251
$|X| \to \mathbf{Z}$ 由下式定义：
$$
x \longmapsto \delta(x)\text{，使得 }
\omega_{X, \overline{x}}^\bullet[-\delta(x)]
\text{ 是下列环上的正规化对偶复形：}
\mathcal{O}_{X, \overline{x}}
$$
该函数是 $|X|$ 上的维数函数。
\end{lemma}

\begin{proof}
取概形 $U$ 及满平展态射 $U\to X$。令 $\omega_U^\bullet$ 为 $U$ 上
与 $\omega_X^\bullet|_U$ 对应的对偶复形。若 $u\in U$ 映到
$x\in|X|$，则 $\mathcal{O}_{X,\overline{x}}$ 是
$\mathcal{O}_{U,u}$ 的严格 Hensel 化。由《对偶复形》引理
\ref{dualizing-lemma-flat-unramified} 可知，若 $\omega^\bullet$ 是
$\mathcal{O}_{U,u}$ 的正规化对偶复形，则
$\omega^\bullet\otimes_{\mathcal{O}_{U,u}}
\mathcal{O}_{X,\overline{x}}$ 是 $\mathcal{O}_{X,\overline{x}}$
的正规化对偶复形。因此，《概形的对偶性》引理
\ref{duality-lemma-dimension-function-scheme} 对概形 $U$ 与复形
$\omega_U^\bullet$ 给出的维数函数 $U\to\mathbf{Z}$，等于
$U\to|X|$ 与 $\delta$ 的复合。利用 $|X|$ 中的特化可提升为
$U$ 中的特化，且 $U$ 中的非平凡特化映到 $X$ 中的非平凡特化
（《适度空间》引理 \ref{decent-spaces-lemma-decent-specialization} 及
\ref{decent-spaces-lemma-decent-no-specializations-map-to-same-point}），
一个简单的拓扑论证表明 $\delta$ 是 $|X|$ 上的维数函数。
\end{proof}






\section{推前的右伴随}
\label{spaces-duality-section-twisted-inverse-image}

\noindent
这是《概形的对偶性》第 \ref{duality-section-twisted-inverse-image} 节
的类似版本。

\begin{lemma}
\label{spaces-duality-lemma-twisted-inverse-image}
\begin{reference}
这几乎就是 \cite[Example 4.2]{Neeman-Grothendieck}。
\end{reference}
设 $S$ 是概形，$f:X\to Y$ 是 $S$ 上拟分离且拟紧的代数空间之间
的态射。函子 $Rf_*:D_\QCoh(X)\to D_\QCoh(Y)$ 有右伴随。
\end{lemma}

\begin{proof}
我们验证《导出范畴》命题 \ref{derived-proposition-brown} 的假设，
由此证明右伴随存在。首先，范畴 $D_\QCoh(\mathcal{O}_X)$ 有直和，
见《空间的导出范畴》引理
\ref{spaces-perfect-lemma-quasi-coherence-direct-sums}；范畴
$D_\QCoh(\mathcal{O}_X)$ 还是紧生成的，
见《空间的导出范畴》定理
\ref{spaces-perfect-theorem-bondal-van-den-Bergh}。
由于 $X$ 与 $Y$ 拟紧且拟分离，$f$ 也如此；见《空间的态射》引理
\ref{spaces-morphisms-lemma-compose-after-separated} 及
\ref{spaces-morphisms-lemma-quasi-compact-permanence}。
因而函子 $Rf_*$ 与直和交换，见《空间的导出范畴》引理
\ref{spaces-perfect-lemma-quasi-coherence-pushforward-direct-sums}。
证明完毕。
\end{proof}

\begin{lemma}
\label{spaces-duality-lemma-twisted-inverse-image-bounded-below}
记号与假设同引理 \ref{spaces-duality-lemma-twisted-inverse-image}。设
$a:D_\QCoh(\mathcal{O}_Y)\to D_\QCoh(\mathcal{O}_X)$ 是 $Rf_*$
的右伴随。则 $a$ 把 $D^+_\QCoh(\mathcal{O}_Y)$ 映入
$D^+_\QCoh(\mathcal{O}_X)$。事实上，存在整数 $N$，使得
$i\leq c$ 时 $H^i(K)=0$ 蕴含 $i\leq c-N$ 时 $H^i(a(K))=0$。
\end{lemma}

\begin{proof}
由《空间的导出范畴》引理
\ref{spaces-perfect-lemma-quasi-coherence-direct-image}，函子 $Rf_*$
具有有限上同调维数。换言之，存在整数 $N$，使得若 $i\geq c$ 时
$H^i(L)=0$，则 $i\geq N+c$ 时 $H^i(Rf_*L)=0$。
设 $K\in D^+_\QCoh(\mathcal{O}_Y)$ 满足 $i\leq c$ 时
$H^i(K)=0$。则
$$
\Hom_{D(\mathcal{O}_X)}(\tau_{\leq c - N}a(K), a(K)) =
\Hom_{D(\mathcal{O}_Y)}(Rf_*\tau_{\leq c - N}a(K), K) = 0
$$
，其中最后一个等号由上述结论得到。显然，这蕴含在 $i\leq c-N$ 时
$H^i(a(K))=0$。
\end{proof}

\noindent
设 $S$ 是概形，$f:X\to Y$ 是 $S$ 上拟分离且拟紧的代数空间之间
的态射。以 $a$ 表示
$Rf_*:D_\QCoh(\mathcal{O}_X)\to D_\QCoh(\mathcal{O}_Y)$ 的右伴随。
对每个 $K\in D_\QCoh(\mathcal{O}_Y)$ 与
$L\in D_\QCoh(\mathcal{O}_X)$，得到典范映射
\begin{equation}
\label{spaces-duality-equation-sheafy-trace}
Rf_*R\SheafHom_{\mathcal{O}_X}(L, a(K))
\longrightarrow
R\SheafHom_{\mathcal{O}_Y}(Rf_*L, K)
\end{equation}
具体而言，该映射构造为复合
$$
Rf_*R\SheafHom_{\mathcal{O}_X}(L, a(K)) \to
R\SheafHom_{\mathcal{O}_Y}(Rf_*L, Rf_*a(K)) \to
R\SheafHom_{\mathcal{O}_Y}(Rf_*L, K)
$$
，其中第一支箭头来自《站点上的上同调》评注
\ref{sites-cohomology-remark-projection-formula-for-internal-hom}，
第二支箭头是该伴随的余单位 $Rf_*a(K)\to K$。

\begin{lemma}
\label{spaces-duality-lemma-iso-on-RSheafHom}
设 $S$ 是概形，$f:X\to Y$ 是 $S$ 上拟紧且拟分离的代数空间之间
的态射，$a$ 是
$Rf_*:D_\QCoh(\mathcal{O}_X)\to D_\QCoh(\mathcal{O}_Y)$ 的右伴随。
设 $L\in D_\QCoh(\mathcal{O}_X)$ 且
$K\in D_\QCoh(\mathcal{O}_Y)$。则映射
(\ref{spaces-duality-equation-sheafy-trace})
$$
Rf_*R\SheafHom_{\mathcal{O}_X}(L, a(K))
\longrightarrow
R\SheafHom_{\mathcal{O}_Y}(Rf_*L, K)
$$
在施加《空间的导出范畴》第
\ref{spaces-perfect-section-better-coherator} 节讨论的函子
$DQ_Y:D(\mathcal{O}_Y)\to D_\QCoh(\mathcal{O}_Y)$ 后成为同构。
\end{lemma}

\begin{proof}
由《空间的导出范畴》引理
\ref{spaces-perfect-lemma-better-coherator}，$DQ_Y$ 存在，故该陈述有意义。
由于 $DQ_Y$ 是包含函子
$D_\QCoh(\mathcal{O}_Y)\to D(\mathcal{O}_Y)$ 的右伴随，为证明引理，
只需证明对任意 $M\in D_\QCoh(\mathcal{O}_Y)$，映射
(\ref{spaces-duality-equation-sheafy-trace}) 诱导双射
% P10-E0252
$$
\Hom_Y(M, Rf_*R\SheafHom_{\mathcal{O}_X}(L, a(K)))
\longrightarrow
\Hom_Y(M, R\SheafHom_{\mathcal{O}_Y}(Rf_*L, K))
$$
。为此，使用下列一串等式：
\begin{align*}
\Hom_Y(M, Rf_*R\SheafHom_{\mathcal{O}_X}(L, a(K)))
& =
\Hom_X(Lf^*M, R\SheafHom_{\mathcal{O}_X}(L, a(K))) \\
& =
\Hom_X(Lf^*M \otimes_{\mathcal{O}_X}^\mathbf{L} L, a(K)) \\
& =
\Hom_Y(Rf_*(Lf^*M \otimes_{\mathcal{O}_X}^\mathbf{L} L), K) \\
& =
\Hom_Y(M \otimes_{\mathcal{O}_Y}^\mathbf{L} Rf_*L, K) \\
& =
\Hom_Y(M, R\SheafHom_{\mathcal{O}_Y}(Rf_*L, K))
\end{align*}
第一个等号由《站点上的上同调》引理
\ref{sites-cohomology-lemma-adjoint} 得到。
% P10-E0253
第二个等号由《站点上的上同调》引理
\ref{sites-cohomology-lemma-internal-hom} 得到；第三个等号来自 $a$ 的构造；
第四个等号由《空间的导出范畴》引理
\ref{spaces-perfect-lemma-cohomology-base-change} 得到（这是关键一步）；
第五个等号由《站点上的上同调》引理
\ref{sites-cohomology-lemma-internal-hom} 得到。
\end{proof}

\begin{example}
\label{spaces-duality-example-iso-on-RSheafHom}
若不施加“拟凝聚化子”$DQ_Y$，引理 \ref{spaces-duality-lemma-iso-on-RSheafHom}
的陈述并不成立。见《概形的对偶性》例
\ref{duality-example-iso-on-RSheafHom}。
\end{example}

\begin{remark}
\label{spaces-duality-remark-iso-on-RSheafHom}
在引理 \ref{spaces-duality-lemma-iso-on-RSheafHom} 的情形中，由《空间的导出范畴》
引理 \ref{spaces-perfect-lemma-pushforward-better-coherator} 有
$$
DQ_Y(Rf_*R\SheafHom_{\mathcal{O}_X}(L, a(K))) =
Rf_* DQ_X(R\SheafHom_{\mathcal{O}_X}(L, a(K)))
$$
。因此，若
$R\SheafHom_{\mathcal{O}_X}(L,a(K))\in D_\QCoh(\mathcal{O}_X)$，
便可从箭头左端“消去”$DQ_Y$。另一方面，若已知
$R\SheafHom_{\mathcal{O}_Y}(Rf_*L,K)\in D_\QCoh(\mathcal{O}_Y)$，
便可从箭头右端“消去”$DQ_Y$。若两者均成立，则
(\ref{spaces-duality-equation-sheafy-trace}) 是同构。结合《空间的导出范畴》引理
\ref{spaces-perfect-lemma-quasi-coherence-internal-hom} 可知，在下列任一
情形中，$Rf_*R\SheafHom_{\mathcal{O}_X}(L,a(K))\to
R\SheafHom_{\mathcal{O}_Y}(Rf_*L,K)$ 是同构：
\begin{enumerate}
\item $L$ 与 $Rf_*L$ 都是完美的；或
\item $K$ 下有界，且 $L$ 与 $Rf_*L$ 都是伪凝聚的。
\end{enumerate}
对于 (2)，使用如下事实：若 $K$ 下有界，则 $a(K)$ 下有界；见引理
\ref{spaces-duality-lemma-twisted-inverse-image-bounded-below}。
\end{remark}

\begin{example}
\label{spaces-duality-example-iso-on-RSheafHom-noetherian}
设 $S$ 是概形，$f:X\to Y$ 是 $S$ 上 Noether 代数空间之间的真态射，
$L\in D^-_{\textit{Coh}}(X)$，且
$K\in D^+_{\QCoh}(\mathcal{O}_Y)$。则映射
$Rf_*R\SheafHom_{\mathcal{O}_X}(L,a(K))\to
R\SheafHom_{\mathcal{O}_Y}(Rf_*L,K)$ 是同构。事实上，由《空间的
导出范畴》引理
\ref{spaces-perfect-lemma-identify-pseudo-coherent-noetherian} 及
\ref{spaces-perfect-lemma-direct-image-coherent}，复形 $L$ 与 $Rf_*L$
都是伪凝聚的，故可应用评注 \ref{spaces-duality-remark-iso-on-RSheafHom} 中的讨论。
\end{example}

\begin{lemma}
\label{spaces-duality-lemma-iso-global-hom}
设 $S$ 是概形，$f:X\to Y$ 是 $S$ 上拟分离且拟紧的代数空间之间
的态射。对所有 $L\in D_\QCoh(\mathcal{O}_X)$ 与
$K\in D_\QCoh(\mathcal{O}_Y)$，(\ref{spaces-duality-equation-sheafy-trace}) 诱导
整体导出 Hom 的同构
$R\Hom_X(L,a(K))\to R\Hom_Y(Rf_*L,K)$。
\end{lemma}

\begin{proof}
由构造（《站点上的上同调》第
\ref{sites-cohomology-section-global-RHom} 节），复形
$$
R\Hom_X(L, a(K)) =
R\Gamma(X, R\SheafHom_{\mathcal{O}_X}(L, a(K))) =
R\Gamma(Y, Rf_*R\SheafHom_{\mathcal{O}_X}(L, a(K)))
$$
与
% P10-E0254
$$
R\Hom_Y(Rf_*L, K) = R\Gamma(Y, R\SheafHom_{\mathcal{O}_X}(Rf_*L, a(K)))
$$
。因此，本引理由引理 \ref{spaces-duality-lemma-iso-on-RSheafHom} 推出。
具体而言，若 $D(\mathcal{O}_Y)$ 中的映射 $E\to E'$ 诱导同构
$DQ_Y(E)\to DQ_Y(E')$，则它诱导拟同构
$R\Gamma(Y,E)\to R\Gamma(Y,E')$。事实上，
$H^i(Y, E) = \Ext^i_Y(\mathcal{O}_Y, E) = \Hom(\mathcal{O}_Y[-i], E) =
\Hom(\mathcal{O}_Y[-i], DQ_Y(E))$。这是因为 $\mathcal{O}_Y[-i]$
属于 $D_\QCoh(\mathcal{O}_Y)$，且 $DQ_Y$
是包含函子 $D_\QCoh(\mathcal{O}_Y)\to D(\mathcal{O}_Y)$ 的右伴随。
\end{proof}









\section{推前的右伴随与基变换（一）}
\label{spaces-duality-section-base-change-map}

\noindent
我们来定义推前函子的右伴随之间的基变换映射。设 $S$ 是概形，
考虑笛卡尔图
\begin{equation}
\label{spaces-duality-equation-base-change}
\vcenter{
\xymatrix{
X' \ar[r]_{g'} \ar[d]_{f'} & X \ar[d]^f \\
Y' \ar[r]^g & Y
}
}
\end{equation}
，其中 $Y'$ 与 $X$ 在 $Y$ 上 {\bf Tor 独立}。以
$$
a  : D_\QCoh(\mathcal{O}_Y) \to D_\QCoh(\mathcal{O}_X)
\quad\text{且}\quad
a' : D_\QCoh(\mathcal{O}_{Y'}) \to D_\QCoh(\mathcal{O}_{X'})
$$
表示 $Rf_*$ 与 $Rf'_*$ 的右伴随（引理
\ref{spaces-duality-lemma-twisted-inverse-image}）。《站点上的上同调》评注
\ref{sites-cohomology-remark-base-change} 的基变换映射给出函子变换
$$
Lg^* \circ Rf_* \longrightarrow Rf'_* \circ L(g')^*
$$
于具有拟凝聚上同调的层的导出范畴上。因此，得到反向的右伴随之间
的变换
$$
a \circ Rg_* \longleftarrow Rg'_* \circ a'
$$

\begin{lemma}
\label{spaces-duality-lemma-flat-precompose-pus}
在图 (\ref{spaces-duality-equation-base-change}) 中，映射
$a\circ Rg_*\leftarrow Rg'_*\circ a'$ 是同构。
\end{lemma}

\begin{proof}
由《空间的导出范畴》引理
\ref{spaces-perfect-lemma-compare-base-change}，对
$D_\QCoh(\mathcal{O}_X)$ 中每个 $K$，基变换映射
$Lg^*\circ Rf_*K\to Rf'_*\circ L(g')^*K$ 都是同构（这里使用了
Tor 独立的假设）。因此，伴随函子之间相应的变换也是同构。
\end{proof}

\noindent
于是可以考虑由下列复合给出的函子
$D_\QCoh(\mathcal{O}_Y)\to D_\QCoh(\mathcal{O}_{X'})$ 的态射：
\begin{equation}
\label{spaces-duality-equation-base-change-map}
L(g')^* \circ a \to
L(g')^* \circ a \circ Rg_* \circ Lg^* \leftarrow
L(g')^* \circ Rg'_* \circ a' \circ Lg^* \to a' \circ Lg^*
\end{equation}
第一支箭头来自伴随映射 $\text{id}\to Rg_*Lg^*$，最后一支箭头来自
伴随映射 $L(g')^*Rg'_*\to\text{id}$。为逆转中间的箭头，需要 Tor
独立的假设，见引理 \ref{spaces-duality-lemma-flat-precompose-pus}。或者，利用
$L(g')^*$ 与 $R(g')_*$ 的伴随性，可以把
(\ref{spaces-duality-equation-base-change-map}) 看作自然变换
$$
a \to a \circ Rg_* \circ Lg^* \leftarrow Rg'_* \circ a' \circ Lg^*
$$
% P10-E0255
，其中第二支箭头仍然可逆。若 $M\in D_\QCoh(\mathcal{O}_X)$ 且
$K\in D_\QCoh(\mathcal{O}_Y)$，则该映射在 Yoneda 函子上由下式给出：
\begin{align*}
\Hom_X(M, a(K))
& =
\Hom_Y(Rf_*M, K) \\
& \to
\Hom_Y(Rf_*M, Rg_* Lg^*K) \\
& =
\Hom_{Y'}(Lg^*Rf_*M, Lg^*K) \\
& \leftarrow
\Hom_{Y'}(Rf'_* L(g')^*M, Lg^*K) \\
& =
\Hom_{X'}(L(g')^*M, a'(Lg^*K)) \\
& =
\Hom_X(M, Rg'_*a'(Lg^*K))
\end{align*}
% P10-E0256
（其中向左的箭头由《空间的导出范畴》引理
\ref{spaces-perfect-lemma-compare-base-change}
给出的基变换定理可知是可逆的），这使构造更为明确。

\medskip\noindent
本节先证明，该基变换映射对于方块的纵横叠合满足若干自然相容性；
这类似于通常基变换映射的《站点上的上同调》评注
\ref{sites-cohomology-remark-compose-base-change} 与
\ref{sites-cohomology-remark-compose-base-change-horizontal}。
建议读者初读时跳过本节余下部分。

\begin{lemma}
\label{spaces-duality-lemma-compose-base-change-maps}
设 $S$ 是概形。考虑交换图
$$
\xymatrix{
X' \ar[r]_k \ar[d]_{f'} & X \ar[d]^f \\
Y' \ar[r]^l \ar[d]_{g'} & Y \ar[d]^g \\
Z' \ar[r]^m & Z
}
$$
，其中各对象均为 $S$ 上拟紧且拟分离的代数空间，两个方块均为
笛卡尔的，并且 $f$ 与 $l$、$g$ 与 $m$ 分别 Tor 独立。则两个方块的
映射 (\ref{spaces-duality-equation-base-change-map}) 复合后给出外侧矩形的基变换映射
（精确陈述见证明）。
\end{lemma}

\begin{proof}
由假设可知 $g\circ f$ 与 $m$ Tor 独立（略去细节），故该陈述有意义。
本证明以 $k^*$ 代替 $Lk^*$，以 $f_*$ 代替 $Rf_*$。令 $a,b,c$
分别为引理 \ref{spaces-duality-lemma-twisted-inverse-image} 对 $f,g,g\circ f$ 给出的
右伴随；带撇版本的记号类似。上方方块对应的箭头是复合
$$
\gamma_{top} :
k^* \circ a \to k^* \circ a \circ l_* \circ l^*
\xleftarrow{\xi_{top}} k^* \circ k_* \circ a' \circ l^* \to a' \circ l^*
$$
，其中 $\xi_{top}:k_*\circ a'\to a\circ l_*$ 是同构（因而可逆），
并且是基变换映射 $l^*\circ f_*\to f'_*\circ k^*$ 的“对偶”箭头。
两端的箭头来自典范映射 $1\to l_*\circ l^*$ 与
$k^*\circ k_*\to1$。类似地，对下方方块有
$$
\gamma_{bot} :
l^* \circ b \to l^* \circ b \circ m_* \circ m^*
\xleftarrow{\xi_{bot}} l^* \circ l_* \circ b' \circ m^* \to b' \circ m^*
$$
对于外侧矩形，得到
$$
\gamma_{rect} :
k^* \circ c \to k^* \circ c \circ m_* \circ m^*
\xleftarrow{\xi_{rect}} k^* \circ k_* \circ c' \circ m^* \to c' \circ m^*
$$
由于 $(g\circ f)_*=g_*\circ f_*$，故 $c=a\circ b$；类似地，
$c'=a'\circ b'$。引理的陈述是，$\gamma_{rect}$ 等于复合
$$
k^* \circ c = k^* \circ a \circ b \xrightarrow{\gamma_{top}}
a' \circ l^* \circ b \xrightarrow{\gamma_{bot}}
a' \circ b' \circ m^* = c' \circ m^*
$$
为证明这一点，考察下图：
$$
\xymatrix{
& & k^* \circ a \circ b \ar[d] \ar[lldd] \\
& & k^* \circ a \circ l_* \circ l^* \circ b \ar[ld] \\
k^* \circ a \circ b \circ m_* \circ m^* \ar[r] &
k^* \circ a \circ l_* \circ l^* \circ b \circ m_* \circ m^* &
k^* \circ k_* \circ a' \circ l^* \circ b \ar[u]_{\xi_{top}} \ar[d] \ar[ld] \\
& k^*\circ k_* \circ a' \circ l^* \circ b \circ m_* \circ m^*
\ar[u]_{\xi_{top}} \ar[rd] &
a' \circ l^* \circ b \ar[d] \\
k^* \circ k_* \circ a' \circ b' \circ m^* \ar[uu]_{\xi_{rect}} \ar[ddrr] &
k^*\circ k_* \circ a' \circ l^* \circ l_* \circ b' \circ m^*
\ar[u]_{\xi_{bot}} \ar[l] \ar[dr] &
a' \circ l^* \circ b \circ m_* \circ m^* \\
& & a' \circ l^* \circ l_* \circ b' \circ m^* \ar[u]_{\xi_{bot}} \ar[d] \\
& & a' \circ b' \circ m^*
}
$$
沿右侧向下得到上述复合，沿左侧向下得到 $\gamma_{rect}$。
由《范畴》引理 \ref{categories-lemma-properties-2-cat-cats}，或者更直接地
由《范畴》定义 \ref{categories-definition-horizontal-composition} 之前的
讨论，此图右侧的所有四边形均交换。因此，只需证明下图
$$
\xymatrix{
a \circ l_* \circ l^* \circ b \circ m_* &
a \circ b \circ m_* \ar[l] \\
k_* \circ a' \circ l^* \circ b \circ m_* \ar[u]_{\xi_{top}} & \\
k_* \circ a' \circ l^* \circ l_* \circ b' \ar[u]_{\xi_{bot}} \ar[r] &
k_* \circ a' \circ b' \ar[uu]_{\xi_{rect}}
}
$$
在逆转箭头 $\xi_{top},\xi_{bot},\xi_{rect}$ 后交换（注意，这不同于
要求原图本身交换）。另一方面，下图
$$
\xymatrix{
& a \circ l_* \circ l^* \circ b \circ m_* \\
a \circ l_* \circ l^* \circ l_* \circ b'
\ar[ru]^{\xi_{bot}} & &
k_* \circ a' \circ l^* \circ b \circ m_* \ar[ul]_{\xi_{top}} \\
& k_* \circ a' \circ l^* \circ l_* \circ b'
\ar[ul]^{\xi_{top}} \ar[ur]_{\xi_{bot}}
}
$$
由《范畴》引理 \ref{categories-lemma-properties-2-cat-cats} 可知交换。
又因为下列两图
% P10-E0257
$$
\vcenter{
\xymatrix{
a \circ l_* \circ l^* \circ b \circ m_* &
a \circ b \circ m \ar[l] \\
a \circ l_* \circ l^* \circ l_* \circ b' \ar[u] &
a \circ l_* \circ b' \ar[l] \ar[u]
}
}
\quad\text{且}\quad
\vcenter{
\xymatrix{
a \circ l_* \circ l^* \circ l_* \circ b' \ar[r] & a \circ l_* \circ b' \\
k_* \circ a' \circ l^* \circ l_* \circ b' \ar[u] \ar[r] &
k_* \circ a' \circ b' \ar[u]
}
}
$$
交换（见所引文献），并且复合
$l_*\to l_*\circ l^*\circ l_*\to l_*$ 是恒等映射，所以只需证明
% P10-E0258
$$
k \circ a' \circ b' \xrightarrow{\xi_{bot}} a \circ l_* \circ b
\xrightarrow{\xi_{top}} a \circ b \circ m_*
$$
等于 $\xi_{rect}$（通过恒等 $a\circ b=c$ 与 $a'\circ b'=c'$）。
这正是《站点上的上同调》评注
\ref{sites-cohomology-remark-compose-base-change} 的对偶陈述，证明完毕。
\end{proof}

\begin{lemma}
\label{spaces-duality-lemma-compose-base-change-maps-horizontal}
设 $S$ 是概形。考虑交换图
$$
\xymatrix{
X'' \ar[r]_{g'} \ar[d]_{f''} & X' \ar[r]_g \ar[d]_{f'} & X \ar[d]^f \\
Y'' \ar[r]^{h'} & Y' \ar[r]^h & Y
}
$$
，其中各对象均为 $S$ 上拟紧且拟分离的代数空间，两个方块均为
笛卡尔的，并且 $f$ 与 $h$、$f'$ 与 $h'$ 分别 Tor 独立。则两个方块的
映射 (\ref{spaces-duality-equation-base-change-map}) 复合后给出外侧矩形的基变换映射
（精确陈述见证明）。
\end{lemma}

\begin{proof}
由假设可知 $f$ 与 $h\circ h'$ Tor 独立（略去细节），故该陈述有意义。
本证明以 $g^*$ 代替 $Lg^*$，以 $f_*$ 代替 $Rf_*$。令 $a,a',a''$
分别为引理 \ref{spaces-duality-lemma-twisted-inverse-image} 对 $f,f',f''$ 给出的
右伴随。右侧方块对应的箭头是复合
$$
\gamma_{right} :
g^* \circ a \to g^* \circ a \circ h_* \circ h^*
\xleftarrow{\xi_{right}} g^* \circ g_* \circ a' \circ h^* \to a' \circ h^*
$$
，其中 $\xi_{right}:g_*\circ a'\to a\circ h_*$ 是同构（因而可逆），
并且是基变换映射 $h^*\circ f_*\to f'_*\circ g^*$ 的“对偶”箭头。
两端的箭头来自典范映射 $1\to h_*\circ h^*$ 与
$g^*\circ g_*\to1$。类似地，对左侧方块有
$$
\gamma_{left} :
(g')^* \circ a' \to (g')^* \circ a' \circ (h')_* \circ (h')^*
\xleftarrow{\xi_{left}}
(g')^* \circ (g')_* \circ a'' \circ (h')^* \to a'' \circ (h')^*
$$
对于外侧矩形，得到
$$
\gamma_{rect} :
k^* \circ a \to
k^* \circ a \circ m_* \circ m^* \xleftarrow{\xi_{rect}}
k^* \circ k_* \circ a'' \circ m^* \to
a'' \circ m^*
$$
，其中 $k=g\circ g'$ 且 $m=h\circ h'$。有
$k^*=(g')^*\circ g^*$ 与 $m^*=(h')^*\circ h^*$。
引理的陈述是，$\gamma_{rect}$ 等于复合
$$
k^* \circ a =
(g')^* \circ g^* \circ a \xrightarrow{\gamma_{right}}
(g')^* \circ a' \circ h^* \xrightarrow{\gamma_{left}}
a'' \circ (h')^* \circ h^* = a'' \circ m^*
$$
为证明这一点，考察下图：
$$
\xymatrix{
& (g')^* \circ g^* \circ a \ar[d] \ar[ddl] \\
& (g')^* \circ g^* \circ a \circ h_* \circ h^* \ar[ld] \\
(g')^* \circ g^* \circ a \circ h_* \circ (h')_* \circ (h')^* \circ h^* &
(g')^* \circ g^* \circ g_* \circ a' \circ h^*
\ar[u]_{\xi_{right}} \ar[d] \ar[ld] \\
(g')^* \circ g^* \circ g_* \circ a' \circ (h')_* \circ (h')^* \circ h^*
\ar[u]_{\xi_{right}} \ar[dr] &
(g')^* \circ a' \circ h^* \ar[d] \\
(g')^* \circ g^* \circ g_* \circ (g')_* \circ a'' \circ (h')^* \circ h^*
\ar[u]_{\xi_{left}} \ar[ddr] \ar[dr] &
(g')^* \circ a' \circ (h')_* \circ (h')^* \circ h^* \\
& (g')^*\circ (g')_* \circ a'' \circ (h')^* \circ h^*
\ar[u]_{\xi_{left}} \ar[d] \\
& a'' \circ (h')^* \circ h^*
}
$$
沿右侧向下得到上述复合，沿左侧向下得到 $\gamma_{rect}$。
由《范畴》引理 \ref{categories-lemma-properties-2-cat-cats}，或者更直接地
由《范畴》定义 \ref{categories-definition-horizontal-composition} 之前的
讨论，此图右侧的所有四边形均交换。因此，只需证明
$$
g_* \circ (g')_* \circ a'' \xrightarrow{\xi_{left}}
g_* \circ a' \circ (h')_* \xrightarrow{\xi_{right}}
a \circ h_* \circ (h')_*
$$
等于 $\xi_{rect}$。这正是《上同调》评注
\ref{cohomology-remark-compose-base-change-horizontal} 的对偶陈述，
证明完毕。
\end{proof}

\begin{remark}
\label{spaces-duality-remark-going-around}
设 $S$ 是概形。考虑交换图
$$
\xymatrix{
X'' \ar[r]_{k'} \ar[d]_{f''} & X' \ar[r]_k \ar[d]_{f'} & X \ar[d]^f \\
Y'' \ar[r]^{l'} \ar[d]_{g''} & Y' \ar[r]^l \ar[d]_{g'} & Y \ar[d]^g \\
Z'' \ar[r]^{m'} & Z' \ar[r]^m & Z
}
$$
，其中各对象均为 $S$ 上拟紧且拟分离的代数空间，所有方块均为
笛卡尔的，并且 $(f,l),(g,m),(f',l'),(g',m')$ 都是 Tor 独立的
态射对。令 $a,a',a'',b,b',b''$ 分别为引理
\ref{spaces-duality-lemma-twisted-inverse-image} 对 $f,f',f'',g,g',g''$ 给出的右伴随。
如下把图中的方块标记为 $A,B,C,D$：
$$
\begin{matrix}
A & B \\
C & D
\end{matrix}
$$
于是，各方块的映射 (\ref{spaces-duality-equation-base-change-map}) 如下（这里使用
$k^*=Lk^*$ 等简写）：
$$
\begin{matrix}
\gamma_A : (k')^* \circ a' \to a'' \circ (l')^* &
\gamma_B : k^* \circ a \to a' \circ l^* \\
\gamma_C : (l')^* \circ b' \to b'' \circ (m')^* &
\gamma_D : l^* \circ b \to b' \circ m^*
\end{matrix}
$$
对于 $2\times1$ 与 $1\times2$ 矩形，另有四个基变换映射：
% P10-E0259
$$
\begin{matrix}
\gamma_{A + B} : (k \circ k')^* \circ a \to a'' \circ (l \circ l')^* \\
\gamma_{C + D} : (l \circ l')^* \circ b \to b'' \circ (m \circ m')^* \\
\gamma_{A + C} : (k')^* \circ (a' \circ b') \to (a'' \circ b'') \circ (m')^* \\
\gamma_{A + C} : k^* \circ (a \circ b) \to (a' \circ b') \circ m^*
\end{matrix}
$$
由引理 \ref{spaces-duality-lemma-compose-base-change-maps-horizontal}，有
$$
\gamma_{A + B} = \gamma_A \circ \gamma_B, \quad
\gamma_{C + D} = \gamma_C \circ \gamma_D
$$
；由引理 \ref{spaces-duality-lemma-compose-base-change-maps}，有
$$
\gamma_{A + C} = \gamma_C \circ \gamma_A, \quad
\gamma_{B + D} = \gamma_D \circ \gamma_B
$$
这里，更准确的写法应为
$\gamma_{A + B} = (\gamma_A \star \text{id}_{l^*}) \circ
(\text{id}_{(k')^*} \star \gamma_B)$，记号同《范畴》第
\ref{categories-section-formal-cat-cat} 节；其他情形类似。不过，我们
继续沿用引理 \ref{spaces-duality-lemma-compose-base-change-maps} 与
\ref{spaces-duality-lemma-compose-base-change-maps-horizontal} 的证明中的记号滥用，
省略与恒等变换的 $\star$ 乘积，因为只要已知变换的源与目标，
% P10-E0260
便可判断应补入哪些恒等变换。综上，先验地得到两个变换
$$
(k')^* \circ k^* \circ a \circ b
\longrightarrow
a'' \circ b'' \circ (m')^* \circ m^*
$$
，即
$$
\gamma_C \circ \gamma_A \circ \gamma_D \circ \gamma_B =
\gamma_{A + C} \circ \gamma_{B + D}
$$
以及
$$
\gamma_C \circ \gamma_D \circ \gamma_A \circ \gamma_B =
\gamma_{C + D} \circ \gamma_{A + B}
$$
本评注的要点是说明这两个变换相等。具体而言，只需证明下图交换：
$$
\xymatrix{
(k')^* \circ a' \circ l^* \circ b \ar[r]_{\gamma_D} \ar[d]_{\gamma_A} &
(k')^* \circ a' \circ b' \circ m^* \ar[d]^{\gamma_A} \\
a'' \circ (l')^* \circ l^* \circ b \ar[r]^{\gamma_D} &
a'' \circ (l')^* \circ b' \circ m^*
}
$$
由《范畴》引理 \ref{categories-lemma-properties-2-cat-cats}，或者更直接地
由《范畴》定义 \ref{categories-definition-horizontal-composition} 之前的
讨论，这确实成立。
\end{remark}








\section{推前的右伴随与基变换（二）}
\label{spaces-duality-section-base-change-II}

\noindent
本节证明，第 \ref{spaces-duality-section-base-change-map} 节的基变换映射在某些情形下
是同构。

\begin{lemma}
\label{spaces-duality-lemma-more-base-change}
在图 (\ref{spaces-duality-equation-base-change}) 中，进一步假设
$g:Y'\to Y$ 是仿射概形之间的态射，且 $f:X\to Y$ 是真的。
则基变换映射 (\ref{spaces-duality-equation-base-change-map}) 在下列情形中诱导同构
$$
L(g')^*a(K) \longrightarrow a'(Lg^*K)
$$
\begin{enumerate}
\item 若 $f$ 平坦且有限表示，则对所有
% P10-E0261
$K\in D_\QCoh(\mathcal{O}_X)$ 均成立；
\item 若 $f$ 完美且 $Y$ 是 Noether 的，则对所有
$K\in D_\QCoh(\mathcal{O}_X)$ 均成立；
\item 若 $g$ 具有有限 Tor 维数且 $Y$ 是 Noether 的，则对
$K\in D_\QCoh^+(\mathcal{O}_X)$ 成立。
\end{enumerate}
\end{lemma}

\begin{proof}
写成 $Y=\Spec(A)$ 与 $Y'=\Spec(A')$。作为仿射态射的基变换，$g'$
是仿射的。取 $D_\QCoh(\mathcal{O}_X)$ 的完美生成元 $M$，见
《空间的导出范畴》定理
\ref{spaces-perfect-theorem-bondal-van-den-Bergh}。则 $L(g')^*M$ 是
$D_\QCoh(\mathcal{O}_{X'})$ 的生成元，见《空间的导出范畴》评注
\ref{spaces-perfect-remark-pullback-generator}。因此，只需证明
(\ref{spaces-duality-equation-base-change-map}) 诱导整体 Hom 复形的同构
\begin{equation}
\label{spaces-duality-equation-iso}
R\Hom_{X'}(L(g')^*M, L(g')^*a(K))
\longrightarrow
R\Hom_{X'}(L(g')^*M, a'(Lg^*K))
\end{equation}
；见《站点上的上同调》第
\ref{sites-cohomology-section-global-RHom} 节，因为这将蕴含
$L(g')^*a(K)\to a'(Lg^*K)$ 的锥为零。证明的结构如下：先证明这些
Hom 复形同构，再在证明的最后部分说明该同构由
(\ref{spaces-duality-equation-iso}) 诱导。

\medskip\noindent
先看左端。由于 $M$ 完美，典范映射
$$
R\Hom_X(M, a(K)) \otimes^\mathbf{L}_A A'
\longrightarrow
R\Hom_{X'}(L(g')^*M, L(g')^*a(K))
$$
由《空间的导出范畴》引理
\ref{spaces-perfect-lemma-affine-morphism-and-hom-out-of-perfect}
可知是同构。再结合引理 \ref{spaces-duality-lemma-iso-global-hom} 的同构
$R\Hom_Y(Rf_*M, K) = R\Hom_X(M, a(K))$
，可知左端等于 $R\Hom_Y(Rf_*M,K)\otimes_A^\mathbf{L}A'$。

\medskip\noindent
再看右端。先使用引理 \ref{spaces-duality-lemma-iso-global-hom} 的同构
$$
R\Hom_{X'}(L(g')^*M, a'(Lg^*K)) = R\Hom_{Y'}(Rf'_*L(g')^*M, Lg^*K)
$$
。由于 $f$ 与 $g$ Tor 独立，由《空间的导出范畴》引理
\ref{spaces-perfect-lemma-compare-base-change}，基变换映射
$Lg^*Rf_*M\to Rf'_*L(g')^*M$ 是同构。因此可将右端改写为
$R\Hom_{Y'}(Lg^*Rf_*M,Lg^*K)$。由于 $Y,Y'$ 仿射，且 $K,Rf_*M$
属于 $D_\QCoh(\mathcal{O}_Y)$（《空间的导出范畴》引理
\ref{spaces-perfect-lemma-quasi-coherence-direct-image}），有典范映射
$$
\beta :
R\Hom_Y(Rf_*M, K) \otimes^\mathbf{L}_A A'
\longrightarrow
R\Hom_{Y'}(Lg^*Rf_*M, Lg^*K)
$$
于 $D(A')$ 中。这就是《代数详论》等式
(\ref{more-algebra-equation-base-change-RHom}) 的箭头；这里使用了
《概形的导出范畴》引理 \ref{perfect-lemma-affine-compare-bounded} 与
\ref{perfect-lemma-quasi-coherence-internal-hom} 在代数表述之间来回转换。
\begin{enumerate}
\item 若 $f$ 平坦且有限表示，则由《空间的导出范畴》引理
\ref{spaces-perfect-lemma-flat-proper-perfect-direct-image-general}，复形
$Rf_*M$ 在 $Y$ 上完美；由《代数详论》引理
\ref{more-algebra-lemma-base-change-RHom} 的 (1)，$\beta$ 是同构。
\item 若 $f$ 完美且 $Y$ 是 Noether 的，则由《空间的态射详论》引理
\ref{spaces-more-morphisms-lemma-perfect-proper-perfect-direct-image}，
复形 $Rf_*M$ 在 $Y$ 上完美；同上，$\beta$ 是同构。
\item 若 $g$ 具有有限 Tor 维数且 $Y$ 是 Noether 的，则复形 $Rf_*M$
在 $Y$ 上伪凝聚（《空间的导出范畴》引理
\ref{spaces-perfect-lemma-direct-image-coherent} 与
\ref{spaces-perfect-lemma-identify-pseudo-coherent-noetherian}）；由
《代数详论》引理 \ref{more-algebra-lemma-base-change-RHom} 的 (4)，
$\beta$ 是同构。
\end{enumerate}
因此，得到与上一段相同的结果。

\medskip\noindent
在证明余下部分，我们说明第二、三段给出的
(\ref{spaces-duality-equation-iso}) 左右两端的认同，事实上正由
(\ref{spaces-duality-equation-iso}) 给出。为使公式易于处理，记
$(-,-)_X=R\Hom_X(-,-)$，以 $-\otimes A'$ 代替
% P10-E0262
$-\otimes_A^\mathbf{L}A'$，并简记 $g^*=Lg^*$、$f_*=Rf_*$。
考虑下列交换图：
% P10-E0263
$$
\xymatrix{
((g')^*M, (g')^*a(K))_{X'} \ar[d] &
(M, a(K))_X \otimes A' \ar[l]^-\alpha \ar[d] &
(f_*M, K)_Y \otimes A' \ar@{=}[l] \ar[d] \\
((g')^*M, (g')^*a(g_*g^*K))_{X'} &
(M, a(g_*g^*K))_X \otimes A' \ar[l]^-\alpha &
(f_*M, g_*g^*K)_Y \otimes A' \ar@{=}[l] \ar@/_4pc/[dd]_{\mu'} \\
((g')^*M, (g')^*g'_*a'(g^*K))_{X'} \ar[u] \ar[d] &
(M, g'_*a'(g^*K))_X \otimes A' \ar[u] \ar[l]^-\alpha \ar[ld]^\mu &
(f_*M, K) \otimes A' \ar[d]^\beta \\
((g')^*M, a'(g^*K))_{X'} &
(f'_*(g')^*M, g^*K)_{Y'} \ar@{=}[l] \ar[r] &
(g^*f_*M, g^*K)_{Y'}
}
$$
标有 $\alpha$ 的箭头，是《空间的导出范畴》引理
\ref{spaces-perfect-lemma-affine-morphism-and-hom-out-of-perfect}
对以 $X',X,Y',Y$ 为顶点之图所给出的映射。由于水平箭头对各变量
具有函子性，图的上部交换。中间的竖直箭头来自引理
\ref{spaces-duality-lemma-flat-precompose-pus} 的可逆变换
$g'_*\circ a'\to a\circ g_*$，因而中间方块交换。沿左侧向下正是
(\ref{spaces-duality-equation-iso})。上方的水平箭头给出证明第二段使用的认同；
包括 $\beta$ 在内的下方水平箭头给出证明第三段使用的认同。
给定 $E\in D(A)$、$E'\in D(A')$ 以及 $D(A)$ 中的
$c:E\to E'$，以 $\mu_c:E\otimes A'\to E'$ 表示由 $c$ 以及限制与
基变换的伴随性所诱导的映射。若 $c$ 明确，便写成 $\mu=\mu_c$，
即从记号中略去 $c$。图中的映射 $\mu$ 属于这种形式，其中 $c$
由下列认同给出：
% P10-E0264
$(M, g'_*a(g^*K))_X = ((g')^*M, a'(g^*K))_{X'}$
。由《空间的导出范畴》评注
\ref{spaces-perfect-remark-multiplication-map}，含 $\mu$ 的三角形交换。

\medskip\noindent
注意，图
$$
\xymatrix{
(M, a(g_*g^*K))_X &
(f_*M, g_* g^*K)_Y \ar@{=}[l] &
(g^*f_*M, g^*K)_{Y'} \ar@{=}[l] \\
(M, g'_* a'(g^*K))_X \ar[u] &
((g')^*M, a'(g^*K))_{X'} \ar@{=}[l] &
(f'_*(g')^*M, g^*K)_{Y'} \ar@{=}[l] \ar[u]
}
$$
由变换 $g'_*\circ a'\to a\circ g_*$ 的定义即知交换。令 $\mu'$ 如上，
对应于认同
% P10-E0265
$(f_*M, g_*g^*K)_X = (g^*f_*M, g^*K)_{Y'}$
，则六边形也交换。因此，只需证明 $\beta$ 等于下列映射与 $\mu'$
的复合：
% P10-E0266
$(f_*M, K)_Y \otimes A' \to (f_*M, g_*g^*K)_X \otimes A'$
。为此，只需证明所诱导的两个映射
$(f_*M,K)_Y\to(g^*f_*M,g^*K)_{Y'}$ 相同。换言之，只需证明下图交换：
$$
\xymatrix{
R\Hom_A(E, K) \ar[rr]_{\text{由 }\beta\text{ 诱导}} \ar[rd] & &
R\Hom_{A'}(E \otimes_A^\mathbf{L} A', K \otimes_A^\mathbf{L} A') \\
& R\Hom_A(E, K \otimes_A^\mathbf{L} A') \ar[ru]
}
$$
对所有 $E,K\in D(A)$ 都成立。而《代数详论》第
\ref{more-algebra-section-base-change-RHom} 节正是如此构造 $\beta$ 的，
故证明完毕。
\end{proof}









\section{推前的右伴随与迹映射}
\label{spaces-duality-section-trace}

\noindent
设 $S$ 是概形，$f:X\to Y$ 是 $S$ 上拟紧且拟分离的代数空间之间
的态射。令 $a:D_\QCoh(\mathcal{O}_Y)\to D_\QCoh(\mathcal{O}_X)$
为引理 \ref{spaces-duality-lemma-twisted-inverse-image} 中的右伴随。由《范畴》第
\ref{categories-section-adjoint} 节，得到函子变换
$$
\text{Tr}_f : Rf_* \circ a \longrightarrow \text{id}
$$
对 $K\in D_\QCoh(\mathcal{O}_Y)$，相应的映射
$\text{Tr}_{f,K}:Rf_*a(K)\longrightarrow K$ 有时称为{\it 迹映射}。
它具有如下性质：刻画右伴随的双射
$$
\Hom_X(L, a(K)) \longrightarrow \Hom_Y(Rf_*L, K)
$$
（其中 $L\in D_\QCoh(\mathcal{O}_X)$）由下式给出：
$$
\varphi \longmapsto \text{Tr}_{f, K} \circ Rf_*\varphi
$$
典范映射 (\ref{spaces-duality-equation-sheafy-trace})
$$
Rf_*R\SheafHom_{\mathcal{O}_X}(L, a(K))
\longrightarrow
R\SheafHom_{\mathcal{O}_Y}(Rf_*L, K)
$$
由与 $\text{Tr}_{f,K}$ 复合得到。本节将考察的每个迹映射都是该迹映射
的特殊情形。在讨论这些特殊情形之前，先证明迹映射的形成与基变换
相交换。

\begin{lemma}[迹映射与基变换]
\label{spaces-duality-lemma-trace-map-and-base-change}
设有图 (\ref{spaces-duality-equation-base-change})。则映射
$1 \star \text{Tr}_f : Lg^* \circ Rf_* \circ a \to Lg^*$ 与
$\text{Tr}_{f'} \star 1 : Rf'_* \circ a' \circ Lg^* \to Lg^*$
通过下列基变换映射相一致：
$\beta : Lg^* \circ Rf_* \to Rf'_* \circ L(g')^*$
（《站点上的上同调》评注 \ref{sites-cohomology-remark-base-change}）
以及 $\alpha : L(g')^* \circ a \to a' \circ Lg^*$
（\ref{spaces-duality-equation-base-change-map}）。
更确切地说，下列函子变换之图交换：
$$
\xymatrix{
Lg^* \circ Rf_* \circ a
\ar[d]_{\beta \star 1} \ar[r]_-{1 \star \text{Tr}_f} &
Lg^* \\
Rf'_* \circ L(g')^* \circ a \ar[r]^{1 \star \alpha} &
Rf'_* \circ a' \circ Lg^* \ar[u]_{\text{Tr}_{f'} \star 1}
}
$$
\end{lemma}

\begin{proof}
本证明以 $f_*$ 代替 $Rf_*$，以 $g^*$ 代替 $Lg^*$，并省略与恒等变换
的 $\star$ 乘积，因为只要已知变换的源与目标，便可判断应补入哪些
恒等变换。% P10-E0268
回忆，$\beta:g^*\circ f_*\to f'_*\circ(g')^*$ 是同构，而 $\alpha$
用 $\beta$ 的伴随同构
$\beta^\vee:g'_*\circ a'\to a\circ g_*$ 定义；见引理
\ref{spaces-duality-lemma-flat-precompose-pus} 及其证明。首先注意，引理之图上方的
水平箭头等于复合
$$
g^* \circ f_* \circ a \to
g^* \circ f_* \circ a \circ g_* \circ g^* \to
g^* \circ g_* \circ g^* \to g^*
$$
，其中第一支箭头是 $(g^*,g_*)$ 的单位，第二支箭头是
$\text{Tr}_f$，第三支箭头是 $(g^*,g_*)$ 的余单位。这直接源于如下
事实：单位与余单位的复合 $g^*\to g^*\circ g_*\circ g^*\to g^*$
是恒等变换。考虑图
$$
\xymatrix{
& g^* \circ f_* \circ a \ar[ld]_\beta \ar[d] \ar[r]_{\text{Tr}_f} & g^* \\
f'_* \circ (g')^* \circ a \ar[dr] &
g^* \circ f_* \circ a \circ g_* \circ g^* \ar[d]_\beta \ar[ru] &
g^* \circ f_* \circ g'_* \circ a' \circ g^* \ar[l]_{\beta^\vee} \ar[d]_\beta &
f'_* \circ a' \circ g^* \ar[lu]_{\text{Tr}_{f'}} \\
& f'_* \circ (g')^* \circ a \circ g_* \circ g^* &
f'_* \circ (g')^* \circ g'_* \circ a' \circ g^* \ar[ru] \ar[l]_{\beta^\vee}
}
$$
% P10-E0267
在此图中，由《范畴》引理 \ref{categories-lemma-properties-2-cat-cats}，
或者更直接地由《范畴》定义
\ref{categories-definition-horizontal-composition} 之前的讨论，两个方块
均交换。由上述讨论，三角形交换。由《范畴》引理
\ref{categories-lemma-transformation-between-functors-and-adjoints}
，下方方块
$$
\xymatrix{
g^* \circ f_* \circ g'_* \circ a' \ar[d]_{\beta^\vee} \ar[r]_-\beta &
f'_* \circ (g')^* \circ g'_* \circ a' \ar[d] \\
g^* \circ f_* \circ a \circ g_* \ar[r] &
\text{id}
}
$$
交换，从而大图中的五边形交换。由于 $\beta$ 与 $\beta^\vee$ 都是
同构，并且按定义沿大图外侧得到
% P10-E0269
$\text{Tr}_f\circ\alpha\circ\beta$，引理得证。
\end{proof}

\noindent
设 $S$ 是概形，$f:X\to Y$ 是 $S$ 上拟紧且拟分离的代数空间之间
的态射。令 $a:D_\QCoh(\mathcal{O}_Y)\to D_\QCoh(\mathcal{O}_X)$
为引理 \ref{spaces-duality-lemma-twisted-inverse-image} 中 $Rf_*$ 的右伴随。由
《范畴》第 \ref{categories-section-adjoint} 节，得到函子变换
$$
\eta_f : \text{id} \to  a \circ Rf_*
$$
，称为该伴随的单位。

\begin{lemma}
\label{spaces-duality-lemma-unit-and-base-change}
设有图 (\ref{spaces-duality-equation-base-change})。则映射
$1 \star \eta_f : L(g')^* \to L(g')^* \circ a \circ Rf_*$ 与
$\eta_{f'} \star 1 : L(g')^* \to a' \circ Rf'_* \circ L(g')^*$
通过下列基变换映射相一致：
$\beta : Lg^* \circ Rf_* \to Rf'_* \circ L(g')^*$
（《站点上的上同调》评注 \ref{sites-cohomology-remark-base-change}）
以及 $\alpha : L(g')^* \circ a \to a' \circ Lg^*$
（\ref{spaces-duality-equation-base-change-map}）。
更确切地说，下列函子变换之图交换：
$$
\xymatrix{
L(g')^* \ar[r]_-{1 \star \eta_f} \ar[d]_{\eta_{f'} \star 1} &
L(g')^* \circ a \circ Rf_* \ar[d]^\alpha \\
a' \circ Rf'_* \circ L(g')^* &
a' \circ Lg^* \circ Rf_* \ar[l]_-\beta
}
$$
\end{lemma}

\begin{proof}
本证明与引理 \ref{spaces-duality-lemma-trace-map-and-base-change} 的证明对偶。
本证明以 $f_*$ 代替 $Rf_*$，以 $g^*$ 代替 $Lg^*$，并省略与恒等变换
的 $\star$ 乘积，因为只要已知变换的源与目标，便可判断应补入哪些
恒等变换。% P10-E0268
回忆，$\beta:g^*\circ f_*\to f'_*\circ(g')^*$ 是同构，而 $\alpha$
用 $\beta$ 的伴随同构
$\beta^\vee:g'_*\circ a'\to a\circ g_*$ 定义；见引理
\ref{spaces-duality-lemma-flat-precompose-pus} 及其证明。首先注意，引理之图左侧的
竖直箭头等于复合
$$
(g')^* \to (g')^* \circ g'_* \circ (g')^* \to
(g')^* \circ g'_* \circ a' \circ f'_* \circ (g')^* \to
a' \circ f'_* \circ (g')^*
$$
，其中第一支箭头是 $((g')^*,g'_*)$ 的单位，第二支箭头是
$\eta_{f'}$，第三支箭头是 $((g')^*,g'_*)$ 的余单位。这直接源于如下
事实：单位与余单位的复合
$(g')^*\to(g')^*\circ(g')_*\circ(g')^*\to(g')^*$ 是恒等变换。
考虑图
$$
\xymatrix{
& (g')^* \circ a \circ f_* \ar[r] &
(g')^* \circ a \circ g_* \circ g^* \circ f_*
\ar[ld]_\beta \\
(g')^* \ar[ru]^{\eta_f} \ar[dd]_{\eta_{f'}} \ar[rd] &
(g')^* \circ a \circ g_* \circ f'_* \circ (g')^* &
(g')^* \circ g'_* \circ a' \circ g^* \circ f_*
\ar[u]_{\beta^\vee} \ar[ld]_\beta \ar[d] \\
& (g')^* \circ g'_* \circ a' \circ f'_* \circ (g')^*
\ar[ld] \ar[u]_{\beta^\vee} &
a' \circ g^* \circ f_* \ar[lld]^\beta \\
a' \circ f'_* \circ (g')^*
}
$$
% P10-E0267
在此图中，由《范畴》引理 \ref{categories-lemma-properties-2-cat-cats}，
或者更直接地由《范畴》定义
\ref{categories-definition-horizontal-composition} 之前的讨论，两个方块
均交换。由上述讨论，三角形交换。由《范畴》引理
\ref{categories-lemma-transformation-between-functors-and-adjoints}
的对偶，下方方块
$$
\xymatrix{
\text{id} \ar[r] \ar[d] &
g'_* \circ a' \circ g^* \circ f_* \ar[d]^\beta \\
g'_* \circ a' \circ g^* \circ f_* \ar[r]^{\beta^\vee} &
a \circ g_* \circ f'_* \circ (g')^*
}
$$
交换，从而大图中的五边形交换。由于 $\beta$ 与 $\beta^\vee$ 都是
同构，并且按定义沿大图外侧得到
$\beta\circ\alpha\circ\eta_f$，引理得证。
\end{proof}





\section{推前的右伴随与拉回}
\label{spaces-duality-section-compare-with-pullback}

\noindent
设 $S$ 是概形，$f:X\to Y$ 是 $S$ 上拟紧且拟分离的代数空间之间
的态射。令 $a$ 为引理 \ref{spaces-duality-lemma-twisted-inverse-image} 中推前的右伴随。
对 $K,L\in D_\QCoh(\mathcal{O}_Y)$，有典范映射
$$
Lf^*K \otimes^\mathbf{L}_{\mathcal{O}_X} a(L)
\longrightarrow
a(K \otimes_{\mathcal{O}_Y}^\mathbf{L} L)
$$
具体而言，该映射伴随于映射
$$
Rf_*(Lf^*K \otimes^\mathbf{L}_{\mathcal{O}_X} a(L)) =
K \otimes^\mathbf{L}_{\mathcal{O}_Y} Rf_*(a(L))
\longrightarrow
K \otimes^\mathbf{L}_{\mathcal{O}_Y} L
$$
（等号由《空间的导出范畴》引理
\ref{spaces-perfect-lemma-cohomology-base-change} 给出），其中使用迹映射
$Rf_*a(L)\to L$。当 $L=\mathcal{O}_Y$ 时，得到映射
\begin{equation}
\label{spaces-duality-equation-compare-with-pullback}
Lf^*K \otimes^\mathbf{L}_{\mathcal{O}_X} a(\mathcal{O}_Y) \longrightarrow a(K)
\end{equation}
；它对 $K$ 具有函子性，并与特异三角相容。

\begin{lemma}
\label{spaces-duality-lemma-compare-with-pullback-perfect}
设 $S$ 是概形，$f:X\to Y$ 是 $S$ 上拟紧且拟分离的代数空间之间
的态射。上面为 $K,L\in D_\QCoh(\mathcal{O}_Y)$ 定义的映射
$Lf^*K \otimes^\mathbf{L}_{\mathcal{O}_X} a(L) \to
a(K \otimes_{\mathcal{O}_Y}^\mathbf{L} L)$
在 $K$ 完美时是同构。特别地，若 $K$ 完美，则
(\ref{spaces-duality-equation-compare-with-pullback}) 是同构。
\end{lemma}

\begin{proof}
令 $K^\vee$ 为 $K$ 的“对偶”，见《站点上的上同调》引理
\ref{sites-cohomology-lemma-dual-perfect-complex}。对
$M\in D_\QCoh(\mathcal{O}_X)$，有
\begin{align*}
\Hom_{D(\mathcal{O}_Y)}(Rf_*M, K \otimes^\mathbf{L}_{\mathcal{O}_Y} L)
& =
\Hom_{D(\mathcal{O}_Y)}(
Rf_*M \otimes^\mathbf{L}_{\mathcal{O}_Y} K^\vee, L) \\
& =
\Hom_{D(\mathcal{O}_X)}(
M \otimes^\mathbf{L}_{\mathcal{O}_X} Lf^*K^\vee, a(L)) \\
& =
\Hom_{D(\mathcal{O}_X)}(M,
Lf^*K \otimes^\mathbf{L}_{\mathcal{O}_X} a(L))
\end{align*}
% P10-E0270
第二个等号由 $a$ 的定义与投影公式（《站点上的上同调》引理
\ref{sites-cohomology-lemma-projection-formula}）得到，也可使用更一般的
《空间的导出范畴》引理
\ref{spaces-perfect-lemma-cohomology-base-change}。结论遂由 Yoneda 引理
推出。
\end{proof}

\begin{lemma}
\label{spaces-duality-lemma-restriction-compare-with-pullback}
设有图 (\ref{spaces-duality-equation-base-change})，并设
$K\in D_\QCoh(\mathcal{O}_Y)$。则下图交换：
$$
\xymatrix{
L(g')^*(Lf^*K \otimes^\mathbf{L}_{\mathcal{O}_X} a(\mathcal{O}_Y))
\ar[r] \ar[d] & L(g')^*a(K) \ar[d] \\
L(f')^*Lg^*K \otimes_{\mathcal{O}_{X'}}^\mathbf{L} a'(\mathcal{O}_{Y'})
\ar[r] & a'(Lg^*K)
}
$$
其中水平箭头分别是 $K$ 与 $Lg^*K$ 的映射
(\ref{spaces-duality-equation-compare-with-pullback})，竖直箭头则使用《站点上的上同调》
评注 \ref{sites-cohomology-remark-base-change} 及
(\ref{spaces-duality-equation-base-change-map}) 构造。
\end{lemma}

\begin{proof}
本证明以 $f_*$ 代替 $Rf_*$，以 $f^*$ 代替 $Lf^*$ 等，并以
$\otimes$ 代替 $\otimes^\mathbf{L}_{\mathcal{O}_X}$ 等。把
(\ref{spaces-duality-equation-compare-with-pullback}) 写成复合
\begin{align*}
f^*K \otimes a(\mathcal{O}_Y)
& \to
a(f_*(f^*K \otimes a(\mathcal{O}_Y))) \\
& \leftarrow
% P10-E0271
a(K \otimes f_*a(\mathcal{O}_K)) \\
& \to
a(K \otimes \mathcal{O}_Y) \\
& \to
a(K)
\end{align*}
。这里，第一支箭头是单位 $\eta_f$；第二支箭头是把 $a$ 施加于
《站点上的上同调》等式
(\ref{sites-cohomology-equation-projection-formula-map}) 所得，而由
《空间的导出范畴》引理
\ref{spaces-perfect-lemma-cohomology-base-change}，该箭头是同构；
第三支箭头是把 $a$ 施加于 $\text{id}_K\otimes\text{Tr}_f$ 所得；
第四支箭头是把 $a$ 施加于同构 $K\otimes\mathcal{O}_Y=K$ 所得。
为证明引理，只需说明这些映射中的每一个都产生引理陈述中的交换方块。
对 $\eta_f$ 与 $\text{Tr}_f$，这分别由引理
\ref{spaces-duality-lemma-unit-and-base-change} 与
\ref{spaces-duality-lemma-trace-map-and-base-change} 给出。对使用《站点上的上同调》等式
(\ref{sites-cohomology-equation-projection-formula-map}) 的箭头，这由
《站点上的上同调》评注
\ref{sites-cohomology-remark-compatible-with-diagram} 给出。
对乘法映射则显然。证明完毕。
\end{proof}











\section{真平坦态射推前的右伴随}
\label{spaces-duality-section-proper-flat}

\noindent
对于拟紧且拟分离的代数空间之间真、平坦且有限表示的态射，推前的
右伴随具有若干显著性质。

\begin{lemma}
\label{spaces-duality-lemma-proper-flat}
设 $S$ 是概形，$Y$ 是 $S$ 上拟紧且拟分离的代数空间，
$f:X\to Y$ 是真、平坦且有限表示的代数空间态射。令 $a$ 为引理
\ref{spaces-duality-lemma-twisted-inverse-image} 中
$Rf_*:D_\QCoh(\mathcal{O}_X)\to D_\QCoh(\mathcal{O}_Y)$ 的右伴随。
则 $a$ 与直和交换。
\end{lemma}

\begin{proof}
设 $P$ 是 $D(\mathcal{O}_X)$ 的完美对象。由《空间的导出范畴》引理
\ref{spaces-perfect-lemma-flat-proper-perfect-direct-image-general}，
复形 $Rf_*P$ 在 $Y$ 上完美。设 $K_i$ 是
$D_\QCoh(\mathcal{O}_Y)$ 中的一族对象。则
\begin{align*}
\Hom_{D(\mathcal{O}_X)}(P, a(\bigoplus K_i))
& =
\Hom_{D(\mathcal{O}_Y)}(Rf_*P, \bigoplus K_i) \\
& =
\bigoplus \Hom_{D(\mathcal{O}_Y)}(Rf_*P, K_i) \\
& =
\bigoplus \Hom_{D(\mathcal{O}_X)}(P, a(K_i))
\end{align*}
，因为完美对象是紧的（《空间的导出范畴》命题
\ref{spaces-perfect-proposition-compact-is-perfect}）。由于
$D_\QCoh(\mathcal{O}_X)$ 有完美生成元（《空间的导出范畴》定理
\ref{spaces-perfect-theorem-bondal-van-den-Bergh}），可知映射
$\bigoplus a(K_i)\to a(\bigoplus K_i)$ 是同构，即 $a$ 与直和交换。
\end{proof}

\begin{lemma}
\label{spaces-duality-lemma-compare-with-pullback-flat-proper}
设 $S$ 是概形，$Y$ 是 $S$ 上拟紧且拟分离的代数空间，
$f:X\to Y$ 是真、平坦且有限表示的代数空间态射。对
$D_\QCoh(\mathcal{O}_Y)$ 的每个对象 $K$，映射
(\ref{spaces-duality-equation-compare-with-pullback}) 都是同构。
\end{lemma}

\begin{proof}
由引理 \ref{spaces-duality-lemma-proper-flat}，$a$ 与直和交换。因此，使
(\ref{spaces-duality-equation-compare-with-pullback}) 成为同构的
$D_\QCoh(\mathcal{O}_Y)$ 对象，构成 $D_\QCoh(\mathcal{O}_Y)$ 的一个
严格满、饱和的三角子范畴，
并且在取直和下保持。由于 $D_\QCoh(\mathcal{O}_Y)$ 是由单个完美对象
生成的模范畴（《空间的导出范畴》定理
\ref{spaces-perfect-theorem-DQCoh-is-Ddga} 与
\ref{spaces-perfect-theorem-bondal-van-den-Bergh}），可如《代数详论》
评注 \ref{more-algebra-remark-P-resolution} 那样论证，只需对一个完美
对象证明 (\ref{spaces-duality-equation-compare-with-pullback}) 是同构。然而，对完美
对象该结论已经成立，见引理
\ref{spaces-duality-lemma-compare-with-pullback-perfect}。
\end{proof}

\begin{lemma}
\label{spaces-duality-lemma-properties-relative-dualizing}
设 $Y$ 是仿射概形，$f:X\to Y$ 是真、平坦且有限表示的代数空间态射。
令 $a$ 为引理 \ref{spaces-duality-lemma-twisted-inverse-image} 中
$Rf_*:D_\QCoh(\mathcal{O}_X)\to D_\QCoh(\mathcal{O}_Y)$ 的右伴随。
则
\begin{enumerate}
\item $a(\mathcal{O}_Y)$ 是 $D(\mathcal{O}_X)$ 中的 $Y$-完美对象；
\item $Rf_*a(\mathcal{O}_Y)$ 的正次数上同调层均为零；
\item $\mathcal{O}_X \to
R\SheafHom_{\mathcal{O}_X}(a(\mathcal{O}_Y), a(\mathcal{O}_Y))$
是同构。
\end{enumerate}
\end{lemma}

\begin{proof}
% P10-E0272
对 $D(\mathcal{O}_X)$ 的完美对象 $E$，有
\begin{align*}
Rf_*(E \otimes_{\mathcal{O}_X}^\mathbf{L} \omega_{X/Y}^\bullet)
& =
Rf_*R\SheafHom_{\mathcal{O}_X}(E^\vee, \omega_{X/Y}^\bullet) \\
& =
R\SheafHom_{\mathcal{O}_Y}(Rf_*E^\vee, \mathcal{O}_Y) \\
& =
(Rf_*E^\vee)^\vee
\end{align*}
第一个等号见《站点上的上同调》引理
\ref{sites-cohomology-lemma-dual-perfect-complex}。第二个等号见引理
\ref{spaces-duality-lemma-iso-on-RSheafHom}、评注 \ref{spaces-duality-remark-iso-on-RSheafHom} 及
《空间的导出范畴》引理
\ref{spaces-perfect-lemma-flat-proper-perfect-direct-image-general}。
第三个等号是对偶的定义。特别地，这些引文还表明所得对象是
$D(\mathcal{O}_Y)$ 的完美对象。于是由《空间的态射详论》引理
\ref{spaces-more-morphisms-lemma-characterize-relatively-perfect}，
$\omega_{X/Y}^\bullet$ 是 $Y$-完美的。这证明了 (1)。

\medskip\noindent
设 $M$ 是 $D_\QCoh(\mathcal{O}_Y)$ 的对象。则
\begin{align*}
\Hom_Y(M, Rf_*a(\mathcal{O}_Y)) & =
\Hom_X(Lf^*M, a(\mathcal{O}_Y)) \\
& =
\Hom_Y(Rf_*Lf^*M, \mathcal{O}_Y) \\
& =
% P10-E0274
\Hom_Y(M \otimes_{\mathcal{O}_Y}^\mathbf{L} Rf_*\mathcal{O}_Y, \mathcal{O}_Y)
\end{align*}
第一个等号由《站点上的上同调》引理
\ref{sites-cohomology-lemma-adjoint} 得到。
% P10-E0273
第二个等号由 $a$ 的构造得到；第三个等号由《空间的导出范畴》引理
\ref{spaces-perfect-lemma-cohomology-base-change} 得到。回忆：对某个 $N$，
$Rf_*\mathcal{O}_X$ 是 Tor 振幅位于 $[0,N]$ 的完美复形；见
《空间的导出范畴》引理
\ref{spaces-perfect-lemma-flat-proper-perfect-direct-image-general}。
因此，可以用位于次数 $[0,N]$ 的有限投射模复形表示
$Rf_*\mathcal{O}_X$（使用《代数详论》引理
\ref{more-algebra-lemma-perfect} 以及 $Y$ 仿射这一事实）。故若对某个
$i>0$ 有 $M=\mathcal{O}_Y[-i]$，则最后一个群为零。由于 $Y$ 仿射，
可知当 $i>0$ 时 $H^i(Rf_*a(\mathcal{O}_Y))=0$。这证明了 (2)。

\medskip\noindent
设 $E$ 是 $D_\QCoh(\mathcal{O}_X)$ 的完美对象。则
\begin{align*}
% P10-E0275
\Hom_X(E, R\SheafHom_{\mathcal{O}_X}(a(\mathcal{O}_Y), a(\mathcal{O}_Y))
& =
\Hom_X(E \otimes_{\mathcal{O}_X}^\mathbf{L} a(\mathcal{O}_Y),
a(\mathcal{O}_Y)) \\
& =
\Hom_Y(Rf_*(E \otimes_{\mathcal{O}_X}^\mathbf{L} a(\mathcal{O}_Y)),
\mathcal{O}_Y) \\
& =
\Hom_Y(Rf_*(R\SheafHom_{\mathcal{O}_X}(E^\vee, a(\mathcal{O}_Y))),
\mathcal{O}_Y) \\
& =
\Hom_Y(R\SheafHom_{\mathcal{O}_Y}(Rf_*E^\vee, \mathcal{O}_Y),
\mathcal{O}_Y) \\
& =
R\Gamma(Y, Rf_*E^\vee) \\
& =
\Hom_X(E, \mathcal{O}_X)
\end{align*}
第一个等号由《站点上的上同调》引理
\ref{sites-cohomology-lemma-internal-hom} 得到。第二个等号是 $a$ 的定义。
第三个等号来自对偶完美复形 $E^\vee$ 的构造，见《站点上的上同调》
引理 \ref{sites-cohomology-lemma-dual-perfect-complex}。第四个等号来自
证明第一段所得的等式
$Rf_*R\SheafHom_{\mathcal{O}_X}(E^\vee, \omega_{X/Y}^\bullet) =
R\SheafHom_{\mathcal{O}_Y}(Rf_*E^\vee, \mathcal{O}_Y)$
。第五个等号由完美复形的二次对偶性（《站点上的上同调》引理
\ref{sites-cohomology-lemma-dual-perfect-complex}）以及下列事实得到：
% P10-E0276
由《空间的导出范畴》引理
\ref{spaces-perfect-lemma-flat-proper-perfect-direct-image-general}
$Rf_*E$ 是完美的。% P10-E0277
最后一个等号是 $f$ 的 Leray 等式。这一串等式实质上通过 Yoneda 引理
证明了 (3)。具体而言，由《空间的导出范畴》引理
\ref{spaces-perfect-lemma-quasi-coherence-internal-hom}，对象
$R\SheafHom(a(\mathcal{O}_Y),a(\mathcal{O}_Y))$ 属于
$D_\QCoh(\mathcal{O}_X)$。在上式中取 $E=\mathcal{O}_X$，得到映射
$\alpha : \mathcal{O}_X \to
R\SheafHom_{\mathcal{O}_X}(a(\mathcal{O}_Y), a(\mathcal{O}_Y))$
，它对应于
$\text{id}_{\mathcal{O}_X} \in \Hom_X(\mathcal{O}_X, \mathcal{O}_X)$.
。由于上述所有同构对 $E$ 具有函子性，$\alpha$ 的锥是
$D_\QCoh(\mathcal{O}_X)$ 中的对象 $C$，并且对所有完美 $E$ 都有
$\Hom(E,C)=0$。由于完美对象生成该范畴（《空间的导出范畴》定理
\ref{spaces-perfect-theorem-bondal-van-den-Bergh}），可知 $\alpha$
是同构。
\end{proof}










\section{真平坦态射的相对对偶复形}
\label{spaces-duality-section-relative-dualizing-proper-flat}

\noindent
受《概形的对偶性》第 \ref{duality-section-proper-flat}、
\ref{duality-section-relative-dualizing-complexes} 节及本章第
\ref{spaces-duality-section-proper-flat} 节内容的启发，作如下定义。

\begin{definition}
\label{spaces-duality-definition-relative-dualizing-proper-flat}
设 $S$ 是概形，$f:X\to Y$ 是 $S$ 上真、平坦且有限表示的代数空间
态射。$X/Y$ 的{\it 相对对偶复形}是一个二元组
$(\omega_{X/Y}^\bullet,\tau)$，它由 $D(\mathcal{O}_X)$ 中的
$Y$-完美对象 $\omega_{X/Y}^\bullet$ 与映射
$$
\tau : Rf_*\omega_{X/Y}^\bullet \longrightarrow \mathcal{O}_Y
$$
组成，并且对任意笛卡尔方块
$$
\xymatrix{
X' \ar[r]_{g'} \ar[d]_{f'} & X \ar[d]^f \\
Y' \ar[r]^g & Y
}
$$
，若 $Y'$ 是仿射概形，则二元组
$(L(g')^*\omega_{X/Y}^\bullet, Lg^*\tau)$
同构于二元组
$(a'(\mathcal{O}_{Y'}), \text{Tr}_{f', \mathcal{O}_{Y'}})$；该二元组见第
\ref{spaces-duality-section-twisted-inverse-image},
\ref{spaces-duality-section-base-change-map},
\ref{spaces-duality-section-base-change-II},
\ref{spaces-duality-section-trace},
\ref{spaces-duality-section-compare-with-pullback} 与
\ref{spaces-duality-section-proper-flat} 节。
\end{definition}

\noindent
这里需要作几点说明。
\begin{enumerate}
\item 在定义 \ref{spaces-duality-definition-relative-dualizing-proper-flat} 中，可以删去
$\omega_{X/Y}^\bullet$ 是 $Y$-完美的这一假设。事实上，让 $Y'$ 遍历
$Y$ 的一个仿射平展覆盖的成员；由引理
\ref{spaces-duality-lemma-properties-relative-dualizing} 可知，
$\omega_{X/Y}^\bullet$ 在 $X$ 的一个平展覆盖的各成员上的限制都是
$Y$-完美的。这蕴含 $\omega_{X/Y}^\bullet$ 是 $Y$-完美的；见
《空间的态射详论》第
\ref{spaces-more-morphisms-section-relatively-perfect} 节。
\item 考虑相对对偶复形 $(\omega_{X/Y}^\bullet,\tau)$ 与定义
\ref{spaces-duality-definition-relative-dualizing-proper-flat} 中的笛卡尔方块。我们把
同构
$(L(g')^*\omega_{X/Y}^\bullet, Lg^*\tau) \cong
(a'(\mathcal{O}_{Y'}), \text{Tr}_{f', \mathcal{O}_{Y'}})$
的存在理解为：对任意 $M'\in D_\QCoh(\mathcal{O}_{X'})$，映射
$$
\Hom_{X'}(M', L(g')^*\omega_{X/Y}^\bullet)
\longrightarrow
\Hom_{Y'}(Rf'_*M', \mathcal{O}_{Y'}),\quad
\varphi' \longmapsto Lg^*\tau \circ Rf'_*\varphi'
$$
是同构。这由 $a'$ 的定义及第 \ref{spaces-duality-section-trace} 节的讨论得到。特别地，
Yoneda 引理保证该同构唯一。
\item 若 $Y$ 本身仿射，则相对对偶复形
$(\omega_{X/Y}^\bullet,\tau)$ 存在，并典范同构于
$(a(\mathcal{O}_Y),\text{Tr}_{f,\mathcal{O}_Y})$，其中 $a$ 是引理
\ref{spaces-duality-lemma-twisted-inverse-image} 中 $Rf_*$ 的右伴随，$\text{Tr}_f$
如第 \ref{spaces-duality-section-trace} 节所述。事实上，给定定义中的图，由引理
\ref{spaces-duality-lemma-more-base-change} 得到同构
$L(g')^*a(\mathcal{O}_Y)\to a'(\mathcal{O}_{Y'})$；由引理
\ref{spaces-duality-lemma-trace-map-and-base-change}，它与迹映射相容。
\end{enumerate}
这些信息恰好足以把局部给出的相对对偶复形粘合成整体对象。建议读者
跳过以下各引理的证明。

\begin{lemma}
\label{spaces-duality-lemma-relative-dualizing-RHom}
% P10-E0278
设 $S$ 是概形，设 $X\to Y$ 是真、平坦且有限表示的代数空间态射。
若 $(\omega_{X/Y}^\bullet,\tau)$ 是相对对偶复形，则
$\mathcal{O}_X \to
R\SheafHom_{\mathcal{O}_X}(\omega_{X/Y}^\bullet, \omega_{X/Y}^\bullet)$
是同构，并且 $Rf_*\omega_{X/Y}^\bullet$ 的正次数上同调层均为零。
\end{lemma}

\begin{proof}
只需在用一个平展于 $Y$ 的仿射概形作基变换后证明；此时结论由引理
\ref{spaces-duality-lemma-properties-relative-dualizing} 得到。
\end{proof}

\begin{lemma}
\label{spaces-duality-lemma-uniqueness-relative-dualizing}
设 $S$ 是概形，设 $X\to Y$ 是真、平坦且有限表示的代数空间态射。
若 $(\omega_j^\bullet,\tau_j)$（$j=1,2$）是 $X/Y$ 上的两个相对对偶
复形，则存在唯一同构
$(\omega_1^\bullet, \tau_1) \to (\omega_2^\bullet, \tau_2)$。
\end{lemma}

\begin{proof}
取平展态射 $g:Y'\to Y$，其中 $Y'$ 是仿射概形，并记基变换为
$X'=Y'\times_YX$。由定义
\ref{spaces-duality-definition-relative-dualizing-proper-flat} 及其后的讨论，存在唯一同构
$\iota : (\omega_1^\bullet|_{X'}, \tau_1|_{Y'}) \to
(\omega_2^\bullet|_{X'}, \tau_2|_{Y'})$。若 $Y''\to Y'$ 还是仿射
概形间的平展态射，且 $X''=Y''\times_YX$，则 $\iota|_{X''}$ 是唯一同构
$(\omega_1^\bullet|_{X''}, \tau_1|_{Y''}) \to
(\omega_2^\bullet|_{X''}, \tau_2|_{Y''})$（由唯一性）。此外，
$$
\text{Ext}^p_{X'}(\omega_1^\bullet|_{X'}, \omega_2^\bullet|_{X'}) = 0,
\quad p < 0
$$
，因为由引理 \ref{spaces-duality-lemma-relative-dualizing-RHom}，
$\mathcal{O}_{X'} \cong
R\SheafHom_{\mathcal{O}_{X'}}(\omega_1^\bullet|_{X'}, \omega_1^\bullet|_{X'})
\cong
R\SheafHom_{\mathcal{O}_{X'}}(\omega_1^\bullet|_{X'}, \omega_2^\bullet|_{X'})$。

\medskip\noindent
选取平展超覆盖 $b:V\to Y$，使得每个
$V_n=\coprod_{i\in I_n}Y_{n,i}$，其中各 $Y_{n,i}$ 均仿射。这可由
《超覆盖》引理 \ref{hypercovering-lemma-hypercovering-object} 与评注
\ref{hypercovering-remark-take-unions-hypercovering-X} 做到（把前一引理
产生的超覆盖换成每一次数上都是不交并的超覆盖）。记
$X_{n,i}=Y_{n,i}\times_YX$ 与 $U_n=V_n\times_YX$，于是得到平展超覆盖
$a:U\to X$（《超覆盖》引理
\ref{hypercovering-lemma-hypercovering-morphism-sites}），且
$U_n=\coprod X_{n,i}$。对 $a:U\to X$ 以及复形
$\omega_1^\bullet$、$\omega_2^\bullet$，均满足《单纯空间》引理
\ref{spaces-simplicial-lemma-fppf-neg-ext-zero-hom} 的假设。因此得到唯一态射
$\iota : \omega_1^\bullet \to \omega_2^\bullet$
，它在 $X_{0,i}$ 上的限制是唯一同构
$(\omega_1^\bullet|_{X_{0, i}}, \tau_1|_{Y_{0, i}}) \to
(\omega_2^\bullet|_{X_{0, i}}, \tau_2|_{Y_{0, i}})$。% P10-E0280
还需证明下图交换：
% P10-E0279
$$
\xymatrix{
Rf_*\omega_1^\bullet \ar[rd]_{\tau_1} \ar[rr]_{Rf_*\iota} & &
Rf_*\omega_1^\bullet \ar[ld]^{\tau_2} \\
& \mathcal{O}_Y
}
$$
不过，$Rf_*\omega_1^\bullet$ 与 $Rf_*\omega_2^\bullet$ 的正次数上同调层
均为零（引理 \ref{spaces-duality-lemma-relative-dualizing-RHom}），故只需限制到各仿射
概形 $Y_{0,i}$ 后证明交换性；在那里它依构造成立。
\end{proof}

\begin{lemma}
\label{spaces-duality-lemma-covering-enough}
% P10-E0281
设 $S$ 是概形，设 $X\to Y$ 是真、平坦且有限表示的代数空间态射。
设 $(\omega^\bullet,\tau)$ 是由 $D(\mathcal{O}_X)$ 中的 $Y$-完美对象与映射
$\tau : Rf_*\omega^\bullet \to \mathcal{O}_Y$。
组成的二元组。设有笛卡尔图
$$
\xymatrix{
X_i \ar[r]_{g_i'} \ar[d]_{f_i} & X \ar[d]^f \\
Y_i \ar[r]^{g_i} & Y
}
$$
，其中 $Y_i$ 仿射，$\{g_i:Y_i\to Y\}$ 是平展覆盖，并且有定义
\ref{spaces-duality-definition-relative-dualizing-proper-flat} 所述的二元组同构
$(\omega^\bullet|_{X_i},\tau|_{Y_i})
\to(a_i(\mathcal{O}_{Y_i}),\text{Tr}_{f_i,\mathcal{O}_{Y_i}})$。
则 $(\omega^\bullet,\tau)$ 是 $X$ 在 $Y$ 上的相对对偶复形。
\end{lemma}

\begin{proof}
令 $g:Y'\to Y$ 以及 $X',f',g',a'$ 如定义
\ref{spaces-duality-definition-relative-dualizing-proper-flat} 所述。置
$((\omega')^\bullet,\tau')=(L(g')^*\omega^\bullet,Lg^*\tau)$。
可以找到一个由仿射概形组成的有限平展覆盖 $\{Y'_j\to Y'\}$，它加细
$\{Y_i\times_YY'\to Y'\}$（《拓扑》引理
\ref{topologies-lemma-etale-affine}）。因此，对每个 $j$，存在 $i_j$ 以及
$Y$ 上的态射 $k_j:Y'_j\to Y_{i_j}$。考虑纤维积
$$
\xymatrix{
X'_j \ar[r]_{h_j'} \ar[d]_{f'_j} &
X' \ar[d]^{f'} \\
Y'_j \ar[r]^{h_j} & Y'
}
$$
。记 $k'_j:X'_j\to X_{i_j}$ 为诱导态射（即 $k_j$ 沿 $f_{i_j}$ 的基变换）。
经态射 $k'_j$ 把给定同构限制到 $Y'_j$，得到二元组同构
% P10-E0282
$((\omega')^\bullet|_{X'_j}, \tau'|_{Y'_j})
\to (a_j(\mathcal{O}_{Y'_j}), \text{Tr}_{f'_j, \mathcal{O}_{Y'_j}})$。
以 $f':X'\to Y'$ 代替 $f:X\to Y$ 后，便归结为下一段所解决的问题。

\medskip\noindent
设 $Y$ 仿射。问题是证明 $(\omega^\bullet,\tau)$ 同构于
$(\omega_{X/Y}^\bullet,\text{Tr})=
(a(\mathcal{O}_Y),\text{Tr}_{f,\mathcal{O}_Y})$。
可以假设覆盖 $\{Y_i\to Y\}$ 由仿射概形间的一个满平展态射
$\{g:Y'\to Y\}$ 给出。事实上，可以先用一个有限子覆盖替换
$\{g_i:Y_i\to Y\}$，再令
$g=\coprod g_i:Y'=\coprod Y_i\to Y$；略去若干细节。置
$X'=Y'\times_YX$，并令 $f',g'$ 如定义
\ref{spaces-duality-definition-relative-dualizing-proper-flat} 所述。此时给定的信息仅为同构
$$
(\omega^\bullet|_{X'}, \tau|_{Y'}) \to
(a'(\mathcal{O}_{Y'}), \text{Tr}_{f', \mathcal{O}_{Y'}})
$$
。由于 $(\omega_{X/Y}^\bullet,\text{Tr})$ 是相对对偶复形（见定义
\ref{spaces-duality-definition-relative-dualizing-proper-flat} 后的讨论），故有唯一同构
$$
(\omega_{X/Y}^\bullet|_{X'}, \text{Tr}|_{Y'}) \to
(a'(\mathcal{O}_{Y'}), \text{Tr}_{f', \mathcal{O}_{Y'}})
$$
。例如，唯一性由引理 \ref{spaces-duality-lemma-uniqueness-relative-dualizing} 得到。
合并以上两个同构，得到同构
$$
\alpha :
(\omega^\bullet|_{X'}, \tau|_{Y'}) \to
(\omega_{X/Y}^\bullet|_{X'}, \text{Tr}|_{Y'})
$$
。% P10-E0283
置 $Y''=Y'\times_YY'$ 与 $X''=Y''\times_YX$。由唯一性，$\alpha$ 到
$X''$ 的两个拉回必相同。又因为 $X'$ 上
$\omega_{X'/Y'}^\bullet$ 的负次数自 Ext 均为零（引理
\ref{spaces-duality-lemma-relative-dualizing-RHom}），而且沿任一投影
$Y'\times_Y\ldots\times_YY'\to Y'$ 拉回后仍然如此（略去一个小细节；
比较引理 \ref{spaces-duality-lemma-uniqueness-relative-dualizing} 的证明），故由
《单纯空间》引理 \ref{spaces-simplicial-lemma-fppf-neg-ext-zero-hom}，
$\alpha$ 下降为 $X$ 上的同构
$\omega^\bullet\to\omega_{X/Y}^\bullet$。
\end{proof}

\begin{lemma}
\label{spaces-duality-lemma-existence-relative-dualizing}
% P10-E0284
设 $S$ 是概形，设 $X\to Y$ 是真、平坦且有限表示的代数空间态射。
存在相对对偶复形 $(\omega_{X/Y}^\bullet,\tau)$。
\end{lemma}

\begin{proof}
选取平展超覆盖 $b:V\to Y$，使每个
$V_n=\coprod_{i\in I_n}Y_{n,i}$，其中各 $Y_{n,i}$ 均仿射。这可由
《超覆盖》引理 \ref{hypercovering-lemma-hypercovering-object} 与评注
\ref{hypercovering-remark-take-unions-hypercovering-X} 做到（把前一引理
产生的超覆盖换成每一次数上都是不交并的超覆盖）。记
$X_{n,i}=Y_{n,i}\times_YX$ 与 $U_n=V_n\times_YX$，于是得到平展超覆盖
$a:U\to X$（《超覆盖》引理
\ref{hypercovering-lemma-hypercovering-morphism-sites}），且
$U_n=\coprod X_{n,i}$。对每个 $n,i$，$X_{n,i}/Y_{n,i}$ 上都存在相对
对偶复形 $(\omega_{n,i}^\bullet,\tau_{n,i})$；见定义
\ref{spaces-duality-definition-relative-dualizing-proper-flat} 后的讨论。
对 $\varphi:[m]\to[n]$ 与 $i\in I_n$，考虑给定超覆盖结构中的态射
% P10-E0285
$g_{\varphi, i} : Y_{n, i} \to Y_{m, \alpha(\varphi)}$ 与
$g'_{\varphi, i} : X_{n, i} \to X_{m, \alpha(\varphi)}$
（《超覆盖》第 \ref{hypercovering-section-hypercovering-sites} 节）。于是
有唯一的二元组同构% P10-E0286
% P10-E0287
$$
\iota_{n, i, \varphi} :
(L(g'_{n, i})^*\omega_{n, i}^\bullet, Lg_{n, i}^*\tau_{n, i})
\longrightarrow
(\omega_{m, \alpha(\varphi)(i)}^\bullet, \tau_{m, \alpha(\varphi)(i)})
$$
；见定义 \ref{spaces-duality-definition-relative-dualizing-proper-flat} 后的讨论。由引理
\ref{spaces-duality-lemma-relative-dualizing-RHom}，$X_{n,i}$ 上
$\omega_{n,i}^\bullet$ 的负次数自 Ext 均为零。记
$(\omega_n^\bullet,\tau_n)$ 为用 $i\in I_n$ 的各二元组
$(\omega_{n,i}^\bullet,\tau_{n,i})$ 构造出的 $U_n/V_n$ 上二元组。
对 $\varphi:[m]\to[n]$ 与 $i\in I_n$，考虑单纯代数空间 $V$ 与 $U$
的结构态射 $g_\varphi:V_n\to V_m$ 与 $g'_\varphi:U_n\to U_m$。
此时有由各分片上的同构构造出的唯一二元组同构
$$
\iota_\varphi :
(L(g'_\varphi)^*\omega_n^\bullet, Lg_\varphi^*\tau_n)
\longrightarrow
(\omega_m^\bullet, \tau_m)
$$
。唯一性保证这些同构满足《单纯空间》定义
\ref{spaces-simplicial-definition-cartesian-derived-modules} 所表述的传递性
条件。对 $a:U\to X$、复形 $\omega_n^\bullet$ 及同构
$\iota_\varphi$，均满足《单纯空间》引理
\ref{spaces-simplicial-lemma-fppf-glue-neg-ext-zero} 的假设\footnote{%
% P10-E0288
该引理只用到 $\omega_0^\bullet$ 以及两个映射
$\delta_1^1,\delta_0^1:[1]\to[0]$。读者可以跳过所引引理证明的开头
几行，因为这里实际上已经给出了模导出范畴中的单纯系统。}。因此得到
$D_\QCoh(\mathcal{O}_X)$ 中的对象 $\omega^\bullet$，以及与两个同构
$\iota_{\delta^1_0}$、$\iota_{\delta^1_1}$ 相容的同构
$\iota_0 : \omega^\bullet|_{U_0} \to \omega_0^\bullet$
。最后应用《单纯空间》引理
\ref{spaces-simplicial-lemma-fppf-neg-ext-zero-hom}，得到唯一态射
$$
\tau : Rf_*\omega^\bullet \longrightarrow \mathcal{O}_Y
$$
，其在 $V_0$ 上的限制与 $\tau_0$ 一致。略去若干细节；例如，可与引理
\ref{spaces-duality-lemma-uniqueness-relative-dualizing} 证明的末尾比较，便可看出为何所需
的负次数 Ext 消失。% P10-E0289
由引理 \ref{spaces-duality-lemma-covering-enough}，二元组
$(\omega^\bullet,\tau)$ 是相对对偶复形，证明完成。
\end{proof}

\begin{lemma}
\label{spaces-duality-lemma-base-change-relative-dualizing}
设 $S$ 是概形。考虑 $S$ 上代数空间的笛卡尔方块
$$
\xymatrix{
X' \ar[d]_{f'} \ar[r]_{g'} & X \ar[d]^f \\
Y' \ar[r]^g & Y
}
$$
。设 $X\to Y$ 真、平坦且有限表示，并设
$(\omega_{X/Y}^\bullet,\tau)$ 是 $f$ 的相对对偶复形。则
$(L(g')^*\omega_{X/Y}^\bullet,Lg^*\tau)$ 是 $f'$ 的相对对偶复形。
\end{lemma}

\begin{proof}
由《空间的态射详论》引理
\ref{spaces-more-morphisms-lemma-base-change-relatively-perfect}，
$L(g')^*\omega_{X/Y}^\bullet$ 是 $Y'$-完美的。定义
\ref{spaces-duality-definition-relative-dualizing-proper-flat} 中的另一个条件由纤维积的
传递性成立。
\end{proof}







\section{与概形情形的比较}
\label{spaces-duality-section-comparison}

\noindent
本节还应补充大量内容。

\begin{lemma}
\label{spaces-duality-lemma-compare}
设 $S$ 是概形，设 $f:X\to Y$ 是 $S$ 上拟紧拟分离代数空间之间的
态射。设 $X$ 与 $Y$ 均可表，并设概形态射 $f_0:X_0\to Y_0$ 表示
$f$（这一记号别扭，但只是临时使用）。令
$a:D_\QCoh(\mathcal{O}_Y)\to D_\QCoh(\mathcal{O}_X)$ 为引理
\ref{spaces-duality-lemma-twisted-inverse-image} 中 $Rf_*$ 的右伴随；令
$a_0:D_\QCoh(\mathcal{O}_{Y_0})\to D_\QCoh(\mathcal{O}_{X_0})$ 为
《概形的对偶性》引理 \ref{duality-lemma-twisted-inverse-image} 中
$Rf_*$ 的右伴随。则下图交换：
$$
\xymatrix{
D_\QCoh(\mathcal{O}_{X_0})
\ar@{=}[rrrrrr]_{\text{《空间的导出范畴》引理
\ref{spaces-perfect-lemma-derived-quasi-coherent-small-etale-site}}}
& & & & & &
D_\QCoh(\mathcal{O}_X) \\
D_\QCoh(\mathcal{O}_{Y_0}) \ar[u]^{a_0}
\ar@{=}[rrrrrr]^{\text{《空间的导出范畴》引理
\ref{spaces-perfect-lemma-derived-quasi-coherent-small-etale-site}}}
& & & & & &
D_\QCoh(\mathcal{O}_Y) \ar[u]_a
}
$$
\end{lemma}

\begin{proof}
这由伴随的唯一性以及《空间的导出范畴》评注
\ref{spaces-perfect-remark-match-total-direct-images} 中的相容性得到。
\end{proof}
